% p4_full_transcript_with_slides.tex
% 编译方式: xelatex p4_full_transcript_with_slides.tex

\documentclass[12pt,a4paper]{article}

\usepackage{fontspec}
\usepackage{xeCJK}
\setCJKmainfont[BoldFont=PingFang SC Semibold]{PingFang SC}
\setCJKsansfont[BoldFont=PingFang SC Semibold]{PingFang SC}
\setCJKmonofont{PingFang SC}
\setmainfont{Palatino}
\setsansfont{Helvetica Neue}
\setmonofont{Menlo}

\usepackage[top=2.5cm, bottom=2.5cm, left=2.5cm, right=2.5cm]{geometry}
\usepackage{amsmath,amssymb}
\usepackage{graphicx}
\usepackage{longtable,booktabs,array}
\usepackage{enumitem}
\usepackage{xcolor}
\usepackage{hyperref}
\usepackage{caption}
\hypersetup{
  colorlinks=true,
  linkcolor=blue!70!black,
  urlcolor=blue!70!black,
  pdfauthor={Auto Generator},
  pdftitle={P4 Context Engineering - Full Transcript with All Slides},
}

\usepackage{parskip}
\setlength{\parindent}{0pt}
\setlength{\parskip}{6pt plus 2pt minus 1pt}

\usepackage{mdframed}
\newmdenv[
  topline=false, bottomline=false, rightline=false,
  linewidth=3pt, linecolor=blue!40,
  innerleftmargin=12pt, innerrightmargin=12pt,
  innertopmargin=8pt, innerbottommargin=8pt,
  backgroundcolor=blue!3, skipabove=10pt, skipbelow=10pt,
]{myquote}

\usepackage{titlesec}
\titleformat{\section}{\Large\bfseries\sffamily}{}{0em}{}[\vspace{-0.5em}\rule{\textwidth}{0.5pt}]
\setcounter{secnumdepth}{0}

\setlength{\emergencystretch}{3em}
\tolerance=1000

\newcommand{\keyslide}[2]{%
  \begin{center}
    \includegraphics[width=0.92\textwidth]{#1}
    \captionof{figure}{#2}
  \end{center}
  \vspace{10pt}
}

\begin{document}

\begin{titlepage}
  \centering
  \vspace*{3cm}
  {\Huge\bfseries\sffamily 第四集 Context Engineering\par}
  \vspace{1.5cm}
  {\Large 全内容对比版（209张幻灯片与完整逐字稿对照呈现）\par}
  \vspace{2cm}
  \vfill
  {\large \today\par}
\end{titlepage}

\newpage
\tableofcontents
\newpage

\section{第四集 Context Engineering (全内容对比版)}

\begin{myquote}
自动生成的带全部图片的逐字稿对应文档。如果图片有外部链接，也会一并提取。
\end{myquote}


\keyslide{../bilibili_video/slides_p4/slide_0001.jpg}{Slide 1}

好,今天这个第二堂课呢？我们要讲什么是 Context Engineering 这是一个最近和热门的 buzzword 那我们今天呢会介绍 Context Engineering 的概念。告诉你说它跟你熟知的 POM Engineering 有什么差异。

那这个 Context Engineering 呢？是今天在 AI Agent 的时代。能够让 AI Agent 成功运行的一个关键的技术。那如果你对 AI Agent 还没有什么概念的话。你其实可以去看我们在今年年初的时候。


\keyslide{../bilibili_video/slides_p4/slide_0002.jpg}{Slide 2}

机器学习这门课讲的一堂课搞懂 AI Agent 的原理。那你可以选择先看这个影片再来听这门课。你也可以听完这门课再来看一堂课搞懂 AI Agent 那两堂课的内容讲的是不太一样的。我们从不同的观点来讲 AI Agent。


\keyslide{../bilibili_video/slides_p4/slide_0003.jpg}{Slide 3}

那两堂课都听完你会有不一样的收获。好,那什么是 Context Engineering 呢？我们不断反复强调说语言模型就是在做文字接龙。语言模型本质上做的事情就是给他一个输入。这个输入叫做 Prop。

他会去预测接下来应该接哪一个 Token 那假设他接的 Token 不对的话怎么办呢？语言模型可以看作是一个韩式叫做 F 它的输入是 X 它的输出就是 F of X。

假设输出不是我们要的。那这个 F of X 里面有两个部分 F 跟 X 所以要嘛就是去改 F 要嘛就是去改 X 那如果你选择的是改 F 的话。

那这件事情叫做训练。它的英文是 Training 或者是叫做学习。它的英文是 Learning 那如果你选择改变的是语言模型的这个韩式 F 里面的参数。


\keyslide{../bilibili_video/slides_p4/slide_0004.jpg}{Slide 4}

改变它的参数。让它有不一样的行为。那这个叫做 Learning 那我们就进入 Machine Learning 的领域。那这个是我们到第五讲后才会提到的。

那今天我们先假设 F 没什么问题。这个语言模型没什么问题。它已经能够尽最大的能力。根据 15 的 Pump 选择最好的接龙的结果。但这件事不一定是真的。

今天语言模型。当然它的能力还是有极限。还是有很多可以再改进的地方。不过多数状况啊。你其实也已经改不了语言模型。

它们都是线上的一个 Close Source 的 Model 它们都是没有开源的模型。所以就算有问题你也改不了。但 F 改不了的话你可以改什么呢？你可以改输入的 X。

你可以帮 F 准备合适的输入 X 让最后得到的 F 的 X 是我们预期的结果。所以今天这一堂课我们就是专注在。怎么给语言模型合适的输入。让语言模型得到正确的输出。

这件事每个人都可以做。就算你手上没有运算资源。你还是可以把 PMP 修好。让语言模型得到你要的输出。主要强调一下。


\keyslide{../bilibili_video/slides_p4/slide_0005.jpg}{Slide 5}

在这一堂课中没有任何模型被训练。我们唯一训练的只有人类而已。好，那刚才讲说我们要修改语言模型的 PMP 那这个不就是 PMP Engineering 吗？那 Context Engineering 跟 PMP Engineering 有什么不同呢？

其实我认为 Context Engineering 的 PMP Engineering 他们没什么本质上的不同 Context Engineering 就是把 PMP Engineering 取一个新的名字。让你觉得潮一点。就好像说早年有 Neural Network。

后来 Neural Network 这个名字就臭掉了。所以有人就把 Neural Network 改成 Deep Learning 其实这个是新瓶装舊久久。其实一样的东西。但是就潮了起来。

所以 Context Engineering 跟 PMP Engineering 我认为它其实就是同的东西。不同的名称，但是其实他们还是有一点点差异。什么样的差异呢？


\keyslide{../bilibili_video/slides_p4/slide_0006.jpg}{Slide 6}

他们指涉的是一样的概念。都是把语言模型的输入弄好。让语言模型得到我们要的输出。但是他们关注的重点是不同的。怎么说呢？

过去当我们讲 PMP Engineering 的时候。也许你马上浮现在脑海中的是。语言模型的输入可能要有特定的格式。你要用某种特定的格式。用 Jason 的 Format。

或者是不同的段落之间。要用紧字号隔开等等。有一阵子有一堆人在教你。怎么用正确的格式来 PMP 语言模型。但今天语言模型越来越强。

通常这些格式已经没有那么重要。往往人看得懂的。语言模型也就看得懂。另外一方面呢？讲到 PMP Engineering 的时候。

可能会让你联想到一系列的神奇咒语。什么神奇咒语呢？在早年这些语言模型还没有那么厉害的时候。他们的输入输出的关系。有时候非常的无厘头。

所以这个时候。你有可能找出神奇的咒语。让语言模型发挥不可思议的力量。比如说这边这个例子。是做在 GPT 3 上面的。

最早的 Chat GPT 是 GPT 3.5 所以 GPT 3 是 Chat GPT 在更早的版本。多数人都没有用过 GPT 3 为什么呢？因为它就是个垃圾。它没有办法被使用。

它非常难被使用。好，这个 GPT 3 它就会有各式各样奇怪的行为。那你可以找到一些神奇咒语来操控它。让它变成你想要的样子。举例来说。

假设我们现在想要 GPT 3 你问它问题的时候。它回应的长度越长越好。那如果你什么都不做。那它平均回应的长度是18.6个字。

如果你在输入的 PMP 跟它说。回答要越长越好。尽可能越长越好。那它可以从18.6个字。进步到23.76个字。

但是我们发现一个神奇的咒语。这个神奇的咒语就是。你只要在矿培里面加一连串的 Waves 就跟它说 Waves Waves Waves Waves Waves Waves 它突然就暴走。


\keyslide{../bilibili_video/slides_p4/slide_0007.jpg}{Slide 7}

\textbf{幻灯片引用链接:}
\begin{itemize}
  \item \href{https://arxiv.org/abs/2206.03931}{https://arxiv.org/abs/2206.03931}
\end{itemize}


\keyslide{../bilibili_video/slides_p4/slide_0008.jpg}{Slide 8}

\textbf{幻灯片引用链接:}
\begin{itemize}
  \item \href{https://arxiv.org/abs/2205.11916}{https://arxiv.org/abs/2205.11916}
\end{itemize}

它答案就变很长。可以长到34.3个字。比你叫它输入输入的答案越长越好。还要更有效，不过这个只有在早年语言模型输入输出的关系。


\keyslide{../bilibili_video/slides_p4/slide_0009.jpg}{Slide 9}

\textbf{幻灯片引用链接:}
\begin{itemize}
  \item \href{https://arxiv.org/abs/2309.03409}{https://arxiv.org/abs/2309.03409}
\end{itemize}


\keyslide{../bilibili_video/slides_p4/slide_0010.jpg}{Slide 10}

\textbf{幻灯片引用链接:}
\begin{itemize}
  \item \href{https://arxivorg/at}{https://arxivorg/at}
\end{itemize}
\begin{itemize}
  \item \href{https://arxiv.org/abs/2}{https://arxiv.org/abs/2}
\end{itemize}


\keyslide{../bilibili_video/slides_p4/slide_0011.jpg}{Slide 11}

\textbf{幻灯片引用链接:}
\begin{itemize}
  \item \href{https://arxiv.org/abs/2312.16171}{https://arxiv.org/abs/2312.16171}
\end{itemize}
\begin{itemize}
  \item \href{https://arxiv.org/at}{https://arxiv.org/at}
\end{itemize}


\keyslide{../bilibili_video/slides_p4/slide_0012.jpg}{Slide 12}

\textbf{幻灯片引用链接:}
\begin{itemize}
  \item \href{https://minimaxir.com/2024/02/chatgpt-tips-analysis/}{https://minimaxir.com/2024/02/chatgpt-tips-analysis/}
\end{itemize}

还很莫名其妙的时候。这种神奇咒语。比较能发挥作用。这边用的这个方法。是出自我们实验室的一篇论文。

那我把论文的连结放在左下角。那我通常我在引用论文的时候。现在我会直接引用这个论文。放在Archive上的连结。那也许不是所有同学都知道什么是Archive。


\keyslide{../bilibili_video/slides_p4/slide_0013.jpg}{Slide 13}

\textbf{幻灯片引用链接:}
\begin{itemize}
  \item \href{https://youtu.be/A3Yx35KrSNO?si=vLfdNcQOhocAN1uc&t=1035}{https://youtu.be/A3Yx35KrSNO?si=vLfdNcQOhocAN1uc\&t=1035}
\end{itemize}


\keyslide{../bilibili_video/slides_p4/slide_0014.jpg}{Slide 14}

\textbf{幻灯片引用链接:}
\begin{itemize}
  \item \href{https://youtu.be/A3Yx35KrSNO?si=vLfdNcQOhocAN1uc&t=1035}{https://youtu.be/A3Yx35KrSNO?si=vLfdNcQOhocAN1uc\&t=1035}
\end{itemize}

所以我还是用几秒钟跟你讲一下什么是Archive Archive是一个放论文的平台。今天在人工智慧领域。因为技术变化的非常的快。所以如果你有新的点子。

你通常不会去投稿。你通常与其等投稿国际会议。那通常要等六周以后。六个月以后，你才会知道论文有没有被接受。

很多时候会研究人员会选择。直接把研究成果公开在Archive这个平台上面。那在Archive这个平台上面。你可以从它的连结。看出这篇论文是什么时候被放上Archive。

比如左下角这个例子。它最后这个连结的数字是2206 那就代表说是2022年6月。被公开的一篇论文。那2022年。

那个是没有Chain GPT的上古时代。那时候人类还只能如毛饮血而已。所以这个是上古时代做的事情。那后来呢？有各式各样的神奇咒语。

很多你也可能都听过。比如说最知名的就是跟模型说。一步一步的思考 Let's fix step by step 模型的能力就起飞了。

后来有人发现说呢？有人发现说呢？不只要告诉他一步一步的思考。还要告诉他呢？请确保你的答案是正确的。

模型可以得到更好的结果。有人发现说呢？如果你叫模型深呼吸在回答问题。可以得到更好的结果。有人发现说呢？

如果你请勒模型跟他说。这题的答案真的对我非常重要。他的正确率可以提高。有人发现说，如果你跟模型说。

如果你答对的话。我就给你一点小费。哇，他的能力居然也会因此而提升。所以模型可能会因为奇奇怪怪的东西。

而让他的正确率。而让他更能够遵守人类给的指令。比如说这边有一个实验的结果是。有人想要测试说啊。这个语言模型最喜欢什么东西。

之前已经讲过说。如果给他小费的话。他呢会做得更好。那语言模型会不会有别的。更吸引他的东西呢？

所以在这一篇ブログ文章里面。作者就试了说。给他泰勒斯的门票。跟他说，如果你做得好。

就会达成世界和平。如果你做得好。你的母亲就会因你而骄傲。或者是如果你做得好。你就会得到珍爱。

看看这个语言模型。他们测试ChairGPT 最喜欢什么东西。那这边他这个实验啊。是要求语言模型。

用200个字写一个故事。所以在这个图上呢？这个我把我的滑鼠叫出来。所以在这个图上呢？黑色的这个纵轴啊。

代表的是正好200个字。所以比200个字多。比200个字少。都算是做得不好。好。

所以整体而言。这个ChairGPT最喜欢什么呢？他最喜欢世界和平。所以你刚才讲说。你做得好。

就可以达成世界和平。可以让他做得最好。他最不在意什么呢？他最不在意自己的母亲。会不会因为他而骄傲。

不过这也是合理的。因为他没有母亲。所以他不在意自己的母亲。也是一个合情合理的结果。总之。


\keyslide{../bilibili_video/slides_p4/slide_0015.jpg}{Slide 15}

曾经有一度，有各式各样的神奇咒语。可以强化语言模型的能力。但是现在，这些神奇的咒语。

越来越不神奇。在一年半之前的。神城市人工智慧。导论正常课里面。我其实就跟大家分享过说。

其实神奇咒语。已经没有以前那么神奇了。这一页投影片。是旧的，一年半前的投影片。

所以如果你是在 2023年6月的时候。你跟GPT3.5 叫他去解一些数学应用问题。没有神奇咒语。

正确率72\% 有神奇咒语。叫他一步一步思考。正确率变88\% 但是过了大概半年之后。

在2024年的2月。你在做一样的实验。如果没有神奇咒语。正确率是85\% 有神奇咒语。

正确率也只能升到89\% 所以你会发现。神奇咒语的效率。是越来越低的。那这个也是一个。

合理的发展方向。因为模型，本来就应该要使尽全力。做到最好，怎么可以说。

叫你一步一步的思考。才开始思考，怎么可以说，给你小费才好好做呢？没有给你小费。

也应该要好好做。所以这些神奇咒语。就逐渐的没有那么神奇了。那神奇咒语不再神奇以后。现在人们关注的重点变了。

那人们现在新的关注的重点。给了新的名字。叫做Context Engineering 这Context Engineering 它真正想要做的。

就是我们自动化的。来管理语言模型的输入。那我现在已经有很多。能够自动做事情的语言模型。我们用语言模型自己。

来管理自己的输入。那如果你要问我说 Context跟Quamp 这两个字眼有什么不同。那我知道说。

可能有些人会把这两个字眼。还是会做出一些差异。那等一下的课堂上。为了要讲解方便。我们就暂时把Context跟Quamp。

画上等号，所以等一下课堂里。当我讲Context的时候。我指的就是Quamp 我指的就是语言模型的输入。

那么现在要讲的。就是管理语言模型输入的方法。好这个是今天的课程大纲。我们会先讲一个语言模型的Context里面。应该包含什么样的资讯。


\keyslide{../bilibili_video/slides_p4/slide_0016.jpg}{Slide 16}


\keyslide{../bilibili_video/slides_p4/slide_0017.jpg}{Slide 17}

然后第二段会讲说。为什么在AI Agent的时代。我们特别需要Context Engineering 然后最后我们会讲一些Context Engineering的基本操作。好那我们就先从Context里面。

应该要包含什么内容开始讲起。一个完整的Context里面。应该要包含什么样的内容呢？最基本的当然是使用者的Quamp 你对语言模型下的指令。

那Quamp里面可以有更多的内容。除了你的任务以外。比如说假设你今天要写封信。跟老师说你meeting要请假。除了告诉语言模型。


\keyslide{../bilibili_video/slides_p4/slide_0018.jpg}{Slide 18}


\keyslide{../bilibili_video/slides_p4/slide_0019.jpg}{Slide 19}

你要执行的任务以外。你最好还是能够提供一些详细的指引。比如说你告诉语言模型说。开头先道歉，然后要说明迟到的理由。

那迟到的理由可能是身体不适。不要讲一些，我今天就是不想meeting之类的理由。然后再说呢？之后会再跟老师呢更新进度。

那老师很忙，所以不要写太长的信。所以你可以交代语言模型说。要在100字以内。然后老师最近心情不太好。

所以跟语言模型说。讲话要非常严肃。不要嬉皮笑脸等等。除了任务的说明以外。你可以加一些额外的条件。

让语言模型的输出。更接近你想要的。那因为语言模型不会读心术。所以如果有一些东西。你希望语言模型做得更好。

那个条件必须要由你来讲清楚。那很多时候呢？你把前提讲清楚。可以让语言模型做得更好。举例来说。

假设你问语言模型。有人告诉我，要用载具吗？这句话是什么意思。那语言模型就开始想。

载具是什么意思。他说载具呢？可能有两种不同的意思。它是交通工具。交通工具确实也可以叫做载具。


\keyslide{../bilibili_video/slides_p4/slide_0020.jpg}{Slide 20}

\textbf{幻灯片引用链接:}
\begin{itemize}
  \item \href{https://www.cinvoice.tw/}{https://www.cinvoice.tw/}
\end{itemize}


\keyslide{../bilibili_video/slides_p4/slide_0021.jpg}{Slide 21}

\textbf{幻灯片引用链接:}
\begin{itemize}
  \item \href{https://www.cinvoice.tw/}{https://www.cinvoice.tw/}
\end{itemize}

那他可能是指某种电子设备或平台。那有人看到这个答案。就觉得不高兴了。想说怎么会觉得载具是交通工具呢？载具不是都是指手机上那个APP。

你买东西之后要扫的吗？怎么会觉得是交通工具呢？但仔细讲讲，载具真的可以指交通工具。你这边并没有告诉语言模型。

你在什么情境下问这个问题。如果你今天讲得更清楚一点说。在超商结账的时候。店员问我要用载具吗？这句话是什么意思。

这个时候语言模型。就可以直确的回答你说。载具就是指电子发票的储存工具。那在这个情境下。因为你已经告诉他。

现在的情境就是在超商结账。所以语言模型就不会觉得载具。是交通工具，或者再举另外一个例子。在做影像处理的时候。


\keyslide{../bilibili_video/slides_p4/slide_0022.jpg}{Slide 22}


\keyslide{../bilibili_video/slides_p4/slide_0023.jpg}{Slide 23}

其实前提也会影响影像处理的结果。这是一张照片。这河上看起来有一只动物。这个照片是在曼谷拍的。去年我去ACL。

我太太也一起去。我去开会的时候。我太太就自己出去玩。她去运河上坐一个船。然后她就发现运河上有一只动物。

就是照片里的这只动物。这只动物离船非常的近。看起来像是鳄鱼。太太就觉得蛮紧张的。这个鳄鱼离人这么近。

感觉他随时是可以爬上船咬人的。但她发现身边的泰国人。一切都不紧张。而且看起来还挺高兴的样子。旁边的泰国人。

试图跟我太太解释这是什么？但是我太太完全听不懂泰文。所以她不知道那些泰国人想要说什么？那些泰国人也不会讲英文。所以没有办法沟通。

所以我太太就拍了这张照片。然后问我说，你看看这是什么动物。然后想说这不是鳄鱼吗？那我就问ChangVT说。

你觉得这是什么动物 ChangVT果然觉得这是鳄鱼。看起来就像鳄鱼。但我想说，这个泰国曼谷的运河里会有鳄鱼吗？

有鳄鱼会让它离人那么近吗？所以我就转念一想。也许我应该提供ChangVT更多的情报。我就问她说现在我们在泰国的曼谷。那水中的动物是什么？


\keyslide{../bilibili_video/slides_p4/slide_0024.jpg}{Slide 24}

这个时候ChangVT就给了我一个不一样的答案。她说这个水中的动物如果在泰国曼谷的话。它看起来是一只水巨蟹。水巨蟹呢在泰国呢其实是一个吉祥的象征。它没有那么常见到但也不是完全见不到。


\keyslide{../bilibili_video/slides_p4/slide_0025.jpg}{Slide 25}

那通常见到水巨蟹就代表说你会有偏财运。那这就是为什么船上的人看到水巨蟹的时候。感觉都蛮高兴的。所以你告诉ChangVT你在哪里拍的一张照片。其实是可以影响它的答案的。


\keyslide{../bilibili_video/slides_p4/slide_0026.jpg}{Slide 26}


\keyslide{../bilibili_video/slides_p4/slide_0027.jpg}{Slide 27}


\keyslide{../bilibili_video/slides_p4/slide_0028.jpg}{Slide 28}


\keyslide{../bilibili_video/slides_p4/slide_0029.jpg}{Slide 29}

那在你的user prompt里面加上范例。也会影响语言模型的答案。比如说我这边呢？要求ChangVT把以下这篇文章呢？用火新闻来改写。

那ChangVT知道火新闻是什么吗？它有一个它想象中的火新闻。它的输出是长这个样子的。而且啊在这个火新闻啊。不是随便乱取的。

比如说你看台北的台。它都把它放到一个框框里面。比如说台北的北。它都是女字边加一个年。它这个修改是有规律的。

它觉得说把一篇文章里面的每一个字。换成另外一个字就是火新闻。但大家知道说火新闻其实指的是。把一些中文字用它的注音符号来替换。所以这边如果我在要求模型写火新闻的时候。


\keyslide{../bilibili_video/slides_p4/slide_0030.jpg}{Slide 30}

\textbf{幻灯片引用链接:}
\begin{itemize}
  \item \href{https://arxiv.org/abs}{https://arxiv.org/abs}
\end{itemize}

我举一个例子跟它说。所谓的火新闻就是要去冒险的人来找我哦。改写成要去冒险的人来找我哦。这个才是我要的火新闻。那当你有给ChangVT这样的举例的时候。

它的输出就变了。它的输出就变成这样。它完全知道所谓的火新闻就是把一些文字改成注音符号。你跟它说我们的环岛之旅。它就知道改成我们的环岛之旅。

它把一些字改成注音符号。所以如果你可以提供清楚的例子。可以影响模型的能力。那把这些例子放到context里面。可以让模型做得更好。

这个早在GPT3的年代就已经知道了。在GPT3那边paper里面他们就说。假设你要叫模型做翻译。你直接跟它说把英文翻成法文。然后叫它翻这个字。

它不见得能够成功做翻译。那个时候模型还很弱。有时候它甚至不知道你到底要它做什么？它连叫它做翻译这样的指令。有时候都看不太懂。

但是如果你可以提供给它一些范例。跟它说所谓的翻译。就是把这个字变成这个字。这个字变成这个字。这个字变成这个字。

听懂了吗？那这样气死的翻译是什么？它就更有机会可以答对。然后在GPT3里面。把提供一些例子这件事情。

叫做In Context Learning 要注意一下。这边的Learning 要放在双引号里面。为什么Learning要放在双引号里面呢？


\keyslide{../bilibili_video/slides_p4/slide_0031.jpg}{Slide 31}

\textbf{幻灯片引用链接:}
\begin{itemize}
  \item \href{https://arxiv.org/abs}{https://arxiv.org/abs}
\end{itemize}


\keyslide{../bilibili_video/slides_p4/slide_0032.jpg}{Slide 32}

\textbf{幻灯片引用链接:}
\begin{itemize}
  \item \href{https://storage.googleapis.com/deepmind-media/gemini/gemini_v1_5_report.pay}{https://storage.googleapis.com/deepmind-media/gemini/gemini\_v1\_5\_report.pay}
\end{itemize}


\keyslide{../bilibili_video/slides_p4/slide_0033.jpg}{Slide 33}

\textbf{幻灯片引用链接:}
\begin{itemize}
  \item \href{https://storage.googleapis.com/deepmind-media/gemini/gemini_v1_5_report.pey}{https://storage.googleapis.com/deepmind-media/gemini/gemini\_v1\_5\_report.pey}
\end{itemize}

因为这边的Learning 并不是传统意义上的Machine Learning 机器学习的学习。这边的参数并没有任何的改变。当我们给模型一些范例的时候。

我们完全不会改变这个模型内部的参数。我们语言模型代表的那个韩式F 它都是一样的。只是因为输入变了。所以输出才会跟着改变。

那输入范例可以改变。语言模型的输出。这个老早就知道了。它叫In Context Learning 那讲到In Context Learning。

我觉得最经典的例子。就是GEMINI 1.5的In Context Learning的能力 GEMINI 1.5刚出的时候。在他们寄出报告里面。就讲了一个GEMINI In Context Learning的神迹。

这个神迹是这样子的。他们叫GEMINI去翻译一种很少人用的语言。这种语言叫做卡拉蒙语。他们叫ChairGPT 把这个句子从英文翻成卡拉蒙语。

卡拉蒙语的资料非常的稀少。据说在世界上可能只有数千个人在用这个语言。那在网路上据说找不到卡拉蒙语的资料。所以这语言模型很有可能。完全不知道什么是卡拉蒙语。


\keyslide{../bilibili_video/slides_p4/slide_0034.jpg}{Slide 34}

\textbf{幻灯片引用链接:}
\begin{itemize}
  \item \href{https://storage.googleapis.com/deepmind-media/gemini/gemini_v1_5_report.pey}{https://storage.googleapis.com/deepmind-media/gemini/gemini\_v1\_5\_report.pey}
\end{itemize}


\keyslide{../bilibili_video/slides_p4/slide_0035.jpg}{Slide 35}

\textbf{幻灯片引用链接:}
\begin{itemize}
  \item \href{https://storage.googleapis.com/deepmind-media/gemini/gemini_v1_5_report.pdf}{https://storage.googleapis.com/deepmind-media/gemini/gemini\_v1\_5\_report.pdf}
\end{itemize}

所以叫它翻译的时候。它是完全翻不了的。但是如果你给这个语言模型。卡拉蒙语的教科书跟你们字典。他读了教科书读了字典。

再读了你的指令以后。再去做文字揭楼。把教科书跟字典放到模型的Context里面。再叫模型做翻译的时候。这时候神奇的事情是。

模型突然之间就会了卡拉蒙语。这一页头影片是 Jerman来1.5技术报告里面的一个实验。他们说，那我们先来测试各个不同的模型。


\keyslide{../bilibili_video/slides_p4/slide_0036.jpg}{Slide 36}

\textbf{幻灯片引用链接:}
\begin{itemize}
  \item \href{https://storage.googleapis.com/deepmind-media/gemini/gemini_v1_5_report.pdf}{https://storage.googleapis.com/deepmind-media/gemini/gemini\_v1\_5\_report.pdf}
\end{itemize}


\keyslide{../bilibili_video/slides_p4/slide_0037.jpg}{Slide 37}

\textbf{幻灯片引用链接:}
\begin{itemize}
  \item \href{https://storage.googleapis.com/deepmind-media/gemini/gemini_v1_5_report.par}{https://storage.googleapis.com/deepmind-media/gemini/gemini\_v1\_5\_report.par}
\end{itemize}


\keyslide{../bilibili_video/slides_p4/slide_0038.jpg}{Slide 38}

\textbf{幻灯片引用链接:}
\begin{itemize}
  \item \href{https://storage.googleapis.com/deepmind-media/gemini/gemini_v1_5_report.pay}{https://storage.googleapis.com/deepmind-media/gemini/gemini\_v1\_5\_report.pay}
\end{itemize}

把卡拉蒙语转英文。还有英文转卡拉蒙语的能力吧？这边的数字代表人类去评。这个模型翻译的结果到底好不好？满分是六分。

这些模型他们的分数都喊不到一分。代表他们完全没办法。把英文转成卡拉蒙语。或卡拉蒙语转成英文。但是如果你提供给这些模型半本教科书。

或者是Jerman来1.5呢？可以读整本教科书。他这边想要炫耀的。其实Jerman来1.5的技术报告里面。特别举这个实验。

一个他非常想要炫耀的点就是。他们可以输入非常长的context 他们context里面可以放很多东西。当时其他模型都只能放半本教科书。他们可以放整本教科书。

当你给Jerman来整本教科书的时候。他得到的评分是四分。或者是5.5分。满分才六分而已。这个是人类的结果。

找了一些人去学习卡拉蒙语。然后也来做翻译。那还是比机器稍微好一点。但是也没好到多少？那这是比较早年的实验了。

其实在一年半以前 2024年的深层试神工智慧导论。其实就讲过这个实验。那最近有人就分析说。那到底这些模型是从教科书的什么地方。


\keyslide{../bilibili_video/slides_p4/slide_0039.jpg}{Slide 39}

\textbf{幻灯片引用链接:}
\begin{itemize}
  \item \href{https://docs.anthropic.com/en/release-}{https://docs.anthropic.com/en/release-}
\end{itemize}


\keyslide{../bilibili_video/slides_p4/slide_0040.jpg}{Slide 40}

\textbf{幻灯片引用链接:}
\begin{itemize}
  \item \href{https://docs.anthropic.com/en/release-}{https://docs.anthropic.com/en/release-}
\end{itemize}

学到了，但这个学要加一个双引号。他不是真正的学。他并没有改变参数。他到底是从教科书里面的什么地方。

让他有了翻译的能力。那这个研究的结果发现说。真正让模型具有翻译能力的。是教科书里面的例句。那些文法的说明什么？

对模型没什么帮助。他没什么看那些文法的说明。真正有帮助的是。那些教科书里面都会有例句。告诉说这个语言的一个句子。

翻成另外一个语言长什么样子。看起来模型是根据那些例句。让他有了翻译新语言的能力。但这边要再强调一下。我们并没有学任何东西。

在incontext learning里面。并没有任何东西被学习。语言模型的文字揭露能力。是不变的，当你把教科书拿掉。

当你没有给语言模型教科书的时候。他就会回复，他没有办法翻译卡拉蒙语的状态。好刚才讲的是user prompt 那这个语言模型呢？

还要有system prompt 那我们其实上周有讲过system prompt的概念 system prompt是开发语言模型的人。这个开发语言模型平台的人。觉得语言模型每次互动的时候。

都需要有的资讯。那cloud的这个模型呢？他的system prompt是公开的。所以可以在以下这个连接里面。看到cloud的system prompt。

长什么样子，那我们就来看一下cloud 最新的版本 OPUS 4.1的system prompt 长什么样子。

这是他的system prompt的前几句话。你看第一句话。开宗明义，就告诉这个模型说。你叫做cloud。

你是由entropic这个公司所打造的。然后再告诉他。今天是几月几号。不然他不知道今天是几月几号啊。所以你知道说。

为什么原模型知道他自己是谁。为什么原模型知道今天是几月几号。这我们上周讲过。他做文字兼容。怎么可能知道自己是谁。

或今天几月几号。这都是已经写好在system prompt里面的东西了。然后今天一个好的模型啊。一个好的服务啊。他的system prompt可以非常的长。

有太多事情要跟这个语言模型交代了 cloud open4.1啊。他system prompt的长度。居然超过2500个字。里面到底都有些什么呢？

你读了这个system prompt以后。你就会知道说。为什么今天这些语言模型。会有这样的行为。可能有可能是system prompt给他的。

所以cloud的system prompt里面包含。那他是谁，他是cloud 然后呢？他还告诉他使用的说明跟限制。

比如说，如果有人问你cloud的API去哪里用。那你就告诉他。读这个网页的内容。然后呢？


\keyslide{../bilibili_video/slides_p4/slide_0041.jpg}{Slide 41}

\textbf{幻灯片引用链接:}
\begin{itemize}
  \item \href{https://docs.anthropic.com}{https://docs.anthropic.com}
\end{itemize}

会告诉他怎么跟使用者互动。所以cloud的system prompt里面写说。如果使用者觉得不高兴。就叫他去按导占的符号。然后他会有些禁止事项。

比如说他会告诉cloud说。不要告诉使用者。怎么合成化学物质。或怎么做核武器。然后呢？

他还会对回应的风格做一些设定。比如说在system prompt里面交代cloud说。不要回答good question 就开头不要讲好问题。可能是单纯做文字接融的时候。

一个问题后面。太容易接出good question了。所以特别交代一下cloud 不要老是用good question来做回应。然后告诉cloud说。

你的知识只到2025年的1月。所以假设有人问的那个问题。显然是2025年1月后才发生的事 cloud就会回答说我不知道。然后呢？

跟cloud讲一些自己的定位。比如说不要说自己是人类。或者是不要说自己有意识。然后呢？告诉cloud说。

如果人类纠正你。不要马上承认错误。今天语言模型很容易承认错误。只要人类说你错了。往往就。

他就会，不管人类说对不对？他都会附和人类。所以这边特别交代cloud说。如果有人说你错了。

你仔细思考以后再回答。不要随便承认自己的错误。所以这个是一个很复杂的SystemProm 那除了SystemProm以外。这些模型。

还需要有对话的历史记录。这些对话的历史记录。就相当于是语言模型的短期记忆。比如说你如果问ChairGPT 你知道隔壁老王是谁吗？


\keyslide{../bilibili_video/slides_p4/slide_0042.jpg}{Slide 42}


\keyslide{../bilibili_video/slides_p4/slide_0043.jpg}{Slide 43}


\keyslide{../bilibili_video/slides_p4/slide_0044.jpg}{Slide 44}

他知道隔壁老王。他说隔壁老王。就是一个常见的幽默角色。是一个虚构的典型人物。那我现在要教他新的知识。

我告诉他说隔壁老王就姓法。他叫法老王，知道吗？这个ChairGPT就知道了。他说隔壁老王等于法老王。

等于法老王，在同一个对话里面。我再继续问他。那你知道隔壁老王是谁吗？他就知道。

他说隔壁老王。法老王，天下无敌王，天下无敌王是他自己加的。他就知道这么加。

我也没有办法。那为什么他知道隔壁老王就是法老王呢？因为当他讲出隔壁老王。隔壁老王，法老王。

天下无敌王，这个句子的时候。其实是根据前面一长串的对话的历史记录。继续去做文字接龙。前面对话的历史记录都已经告诉他。

隔壁老王就是法老王了。所以他当然知道隔壁老王就是法老王。那很多人在跟ChairGPT对话的时候。往往有一个误解。你以为你跟他对话。

就可以教他新东西。就可以叫他做学习。我只能告诉你。你跟他对话的时候。你是没有办法改变他的参数的。

你没有办法单评跟一个语言模型对话。就改变他的参数。也就是说，你没有办法单评跟一个语言模型对话。就训练他。

改变他真正的内部的行为。所以现在我在一个对话里面。教会ChairGPT说隔壁老王就是法老王。那再开一个新的对话。再问他隔壁老王是谁。

他又变成原来那个样子。他还是不知道隔壁老王是法老王。不过这个现象就是。仅止于2024年9月之前。在2024年9月之后。


\keyslide{../bilibili_video/slides_p4/slide_0045.jpg}{Slide 45}

ChairGPT有了长期的记忆。那等一下在往后的课程里面。我们还会讲一下。这个长期记忆的能力是怎么来的。我们先来看一下有长期记忆以后。


\keyslide{../bilibili_video/slides_p4/slide_0046.jpg}{Slide 46}

这个模型的能力有什么不同。如果你想要开启这个ChairGPT的长期记忆的话。你就在那个自定ChairGPT这个地方。在左下角个人帐篷那边点自定ChairGPT 里面呢在个人化的栏位里面有一个叫记忆的选项。

它其实有两种不同的记忆方式。你把记忆呢开起来。这个模型呢就有了记忆的能力。那虽然等一下在课程里面不会明确的倒数你说 ChairGPT实际上怎么实践它的长期记忆。

不过在课程的结尾会讲一些模型处理记忆的方法。那你可以把模型处理的记忆的方法。跟ChairGPT它的行为对照一下。猜猜看ChairGPT是怎么实践它的长期记忆的。总之啊现在当ChairGPT你问它一个问题的时候。


\keyslide{../bilibili_video/slides_p4/slide_0047.jpg}{Slide 47}


\keyslide{../bilibili_video/slides_p4/slide_0048.jpg}{Slide 48}

在你的问题之前已经被植入了长期记忆。而这些长期记忆是你看不到的。那我觉得这个OpenAI真的是行销鬼才啊。他们帮模型做了长期记忆以后。那怕大家不知道嘛？

因为Rate跟大家说。我这个模型有长期记忆。你会不知道要怎么用这个模型。怎么展现它的长期记忆。他们就做了一个行销的活动。

叫大家去问ChairGPT说我是什么样的人。因为它现在有长期记忆。它会利用你跟它互动的过去的历史记录。在历史记录形成长期记忆。再根据长期记忆来回答问题。

所以当你问它我是什么样的人的时候。它就可以根据它长期记忆里面跟你互动的理解。开始说你是什么样的人。这个ChairGPT说我是什么样的人。基本上它只会说好话。

那有趣的地方是。它知道深层式人工智慧与机器学习。导论这门课，当然因为在过去的历史记录中。跟它解敌过这门课。


\keyslide{../bilibili_video/slides_p4/slide_0049.jpg}{Slide 49}


\keyslide{../bilibili_video/slides_p4/slide_0050.jpg}{Slide 50}


\keyslide{../bilibili_video/slides_p4/slide_0051.jpg}{Slide 51}


\keyslide{../bilibili_video/slides_p4/slide_0052.jpg}{Slide 52}

\textbf{幻灯片引用链接:}
\begin{itemize}
  \item \href{https://www.linkedin.com/posts/petergyang_google-ai-overview-suggests-adding-glue-to-activ'ty}{https://www.linkedin.com/posts/petergyang\_google-ai-overview-suggests-adding-glue-to-activ'ty}
\end{itemize}


\keyslide{../bilibili_video/slides_p4/slide_0053.jpg}{Slide 53}

\textbf{幻灯片引用链接:}
\begin{itemize}
  \item \href{https://www.linkedin.com/posts/petergyang_google-ai-overview-suggests-adding-glue-to-activity}{https://www.linkedin.com/posts/petergyang\_google-ai-overview-suggests-adding-glue-to-activity}
\end{itemize}


\keyslide{../bilibili_video/slides_p4/slide_0054.jpg}{Slide 54}

\textbf{幻灯片引用链接:}
\begin{itemize}
  \item \href{https://www.linkedin.com/posts/petergyang_google-ai-overview-suggests-adding-glue-to-activ'ty}{https://www.linkedin.com/posts/petergyang\_google-ai-overview-suggests-adding-glue-to-activ'ty}
\end{itemize}

所以它知道它在长期记忆里面。有放这门课相关的资讯。那除了记忆以外。我们也会希望可以提供这个语言模型。来自其他资料源的相关资讯。

因为今天语言模型它的知识可能是有限的。而且它的知识可能是过时的。所以我们往往希望能够自动的。提供语言模型一些额外的资讯。怎么样提供语言模型额外的资讯呢？

最常用的做法就是透过搜寻引擎。所以今天你在问一个语言模型任何问题之前。你可以先把这个问题丢到网路上。或某一个你自己的资料库。先做搜寻。

得到额外的资讯。语言模型，根据你给它的quamp 还有额外的资讯。再去做文字接龙。

它就更有可能可以得到正确的答案。那在我们第二讲的作业二。其实就是要叫同学们来实作这个技术。那这个技术有个鼎鼎大名的名字。就是Retrieval Augmented Generation。

也就是RAG 那我们在课堂上。不会讲太多RAG的东西。等一下主教再讲作业二的时候。会非常详细的跟你讲。

RAG是怎么被实践的。好，那今天其实你要在确GPT的平台上。实践RAG是轻而易举了。因为确GPT。

现在你要让它搭配搜寻引擎来使用。你其实在对话框下选个搜寻。它就可以搭配搜寻引擎来使用了。在这边想要提醒大家的是。并不是语言模型结合搜寻引擎以后。

就天下无敌，不要忘了语言模型总是在做文字揭楼。所以就算是搭配了搜寻引擎提供的最新的资讯。它还是有可能会犯错的。一个最经典的。

一直被拿出来津津乐道的例子。就是Google刚在测试这个AI Overview的能力的时候。你知道今天你在Google上输入一个问题。除了搜寻到网页之外。也会得到一个语言模型的答案。

那这就是一个非常经典的RIG的应用。语言模型根据搜寻到的结果。给了你一个他用文字揭楼写出来的答案。在AI Overview还在测试期间的时候。就有人问了一个问题。

我的起司黏不在披萨上。你说要怎么办。这个AI Overview这个语言模型就说。那很简单，我们就把八分之一勺无毒的浇水。

我们就可以把这个起司黏在披萨上了。而且语言模型这个答案。显然不在开玩笑。还有强调要用无毒的浇水。那为什么语言模型会这样回答呢？

可能就是因为在他搜寻到的文章里面。有一个Ready的Po文。有人开玩笑说。只要用八分之一勺的浇水。就可以把起司黏在披萨上了。

语言模型就照单全收。得到错误的答案。其实啊，就算是现在比较新的 Chair GPT-4O。

你开启搜寻引擎的时候。它还是有可能会犯错的。比如说我这边。我问GPT-4O 要强调我有开启搜寻引擎的功能。


\keyslide{../bilibili_video/slides_p4/slide_0055.jpg}{Slide 55}

\textbf{幻灯片引用链接:}
\begin{itemize}
  \item \href{https://chatgpt.com/share/688f1941-68c8-8010-b1f0-5dd30a5af184}{https://chatgpt.com/share/688f1941-68c8-8010-b1f0-5dd30a5af184}
\end{itemize}


\keyslide{../bilibili_video/slides_p4/slide_0056.jpg}{Slide 56}

\textbf{幻灯片引用链接:}
\begin{itemize}
  \item \href{https://chatgpt.com/share/688f1941-68c8-8010-b1f0-5dd30a5af184}{https://chatgpt.com/share/688f1941-68c8-8010-b1f0-5dd30a5af184}
\end{itemize}
\begin{itemize}
  \item \href{https://taicatw.net/fall-114/}{https://taicatw.net/fall-114/}
\end{itemize}
\begin{itemize}
  \item \href{https://taicatwnet/}{https://taicatwnet/}
\end{itemize}


\keyslide{../bilibili_video/slides_p4/slide_0057.jpg}{Slide 57}

我叫他介绍一下。我们开卡，台湾大专院校。人工智慧学成联盟。这个年度上学期呢？

有哪些课程，并提供官网网址。那在搜寻网路之后。在进行回答以后。他的答案当然会正确很多。

上周我们也试过一样的例子。在关闭搜寻引擎的前提之下。模型没办法回答这个问题。他给了一些网址。但网址也是错的。

现在在有搜寻引擎的情况下。那网址就会是正确的。但这并不代表。有了搜寻引擎模型的答案。通通都是正确的。

比如说如果你仔细阅读他的介绍的话。他会说十门课都是镜像课程。但实际上，并不是所有的课程都是镜像课程。那到底什么是镜像课程。

没那么重要，反正确计提也不知道。他不知道从哪里读到所有的课程。都是镜像课程。但这是一个错误的资讯。

所以不要忘了。就算是有搜寻引擎。语言模型是根据搜寻引擎搜寻到的文字。再去做文字接融。他仍然有犯错的可能性。

好,那今天语言模型啊。往往能够使用工具。那使用搜寻引擎也算是使用工具的一种。但除了使用搜寻引擎以外。今天各个知名的语言模型。

Cloud啊,Germini啊,ChairGPT啊。他们都可以搭配很多工具来使用。比如说刚才讲的几个模型。都可以搜寻你的Gmail 他们都可以去阅读你的Google Calendar。

知道你今天要做什么事情。甚至啊,Germini呢？是可以直接去改你的Google Calendar 你可以直接跟他说。我明天上午九点到九点半。


\keyslide{../bilibili_video/slides_p4/slide_0058.jpg}{Slide 58}

要跟王小明meeting 寄在Google日历上。他会真的直接去连到Google日历。真的改你的Google日历。真的加这个项目上去。

所以他可以当做一个助理来使用。好,那模型到底是怎么使用工具的呢? 你需要在你的context里面。先写好使用工具的方法。首先你可能要教他一般的工具。

所有的工具要怎么使用。那这边是一个实际上可以用的例子。那我这边跟他说呢？如果你的知识无法回答问题。那就使用工具。


\keyslide{../bilibili_video/slides_p4/slide_0059.jpg}{Slide 59}


\keyslide{../bilibili_video/slides_p4/slide_0060.jpg}{Slide 60}

你把使用工具的指令。放在Tor跟斜线Tor这两个符号中间。使用完工具后呢？你会看到工具的输出。工具的输出放在Open跟斜线Open中间。


\keyslide{../bilibili_video/slides_p4/slide_0061.jpg}{Slide 61}


\keyslide{../bilibili_video/slides_p4/slide_0062.jpg}{Slide 62}

那使用工具，其实现在有很多不同的方式。有人可能会告诉你说。这个工具呢？的操作说明呢？

要用什么JSON format 机器才看得懂等等。其实不是，今天语言模型很厉害。所以你基本上呢？


\keyslide{../bilibili_video/slides_p4/slide_0063.jpg}{Slide 63}


\keyslide{../bilibili_video/slides_p4/slide_0064.jpg}{Slide 64}


\keyslide{../bilibili_video/slides_p4/slide_0065.jpg}{Slide 65}


\keyslide{../bilibili_video/slides_p4/slide_0066.jpg}{Slide 66}


\keyslide{../bilibili_video/slides_p4/slide_0067.jpg}{Slide 67}


\keyslide{../bilibili_video/slides_p4/slide_0068.jpg}{Slide 68}


\keyslide{../bilibili_video/slides_p4/slide_0069.jpg}{Slide 69}


\keyslide{../bilibili_video/slides_p4/slide_0070.jpg}{Slide 70}

你可以用一般人类看得懂的文字。来说明工具要怎么使用。语言模型可以看懂。他就有办法用这些工具。好那这边呢？

可以告诉他一些特定工具的使用方式。我们这边有个特定的工具。叫做Temperature 他可以查询某个地方。某个时间的温度。

那给他的使用范例。告诉他说Temperature 夸号台北时间。你就可以呼叫这一个韩式。他可以告诉你某一个城市。

某个时间的温度是多少？好那有了使用工具的方式。还有使用工具的范例等等之后。把这些内容加上user prompt 都去丢给语言模型。

所以我们现在的context里面。有包含工具的使用说明。所以你要让模型使用工具的时候。你就要在你的这个context里面。加上工具的使用说明。

好有了这个工具的使用说明以后呢？语言模型就真的有办法使用这些工具了。当你问他一个问题。说2025年某年某月某日。高雄的气温母如何？

这个时候语言模型呢？就会真的接出使用工具的指令。他就会说Tour Temperature 高雄时间斜线托。但是要注意一下。

语言模型的输出就是一串文字。这串文字无法让你真的去使用工具。它就是一串文字。它就是放在那边。就是一串文字。


\keyslide{../bilibili_video/slides_p4/slide_0071.jpg}{Slide 71}


\keyslide{../bilibili_video/slides_p4/slide_0072.jpg}{Slide 72}


\keyslide{../bilibili_video/slides_p4/slide_0073.jpg}{Slide 73}


\keyslide{../bilibili_video/slides_p4/slide_0074.jpg}{Slide 74}


\keyslide{../bilibili_video/slides_p4/slide_0075.jpg}{Slide 75}


\keyslide{../bilibili_video/slides_p4/slide_0076.jpg}{Slide 76}


\keyslide{../bilibili_video/slides_p4/slide_0077.jpg}{Slide 77}


\keyslide{../bilibili_video/slides_p4/slide_0078.jpg}{Slide 78}


\keyslide{../bilibili_video/slides_p4/slide_0079.jpg}{Slide 79}


\keyslide{../bilibili_video/slides_p4/slide_0080.jpg}{Slide 80}


\keyslide{../bilibili_video/slides_p4/slide_0081.jpg}{Slide 81}


\keyslide{../bilibili_video/slides_p4/slide_0082.jpg}{Slide 82}


\keyslide{../bilibili_video/slides_p4/slide_0083.jpg}{Slide 83}


\keyslide{../bilibili_video/slides_p4/slide_0084.jpg}{Slide 84}


\keyslide{../bilibili_video/slides_p4/slide_0085.jpg}{Slide 85}


\keyslide{../bilibili_video/slides_p4/slide_0086.jpg}{Slide 86}


\keyslide{../bilibili_video/slides_p4/slide_0087.jpg}{Slide 87}


\keyslide{../bilibili_video/slides_p4/slide_0088.jpg}{Slide 88}


\keyslide{../bilibili_video/slides_p4/slide_0089.jpg}{Slide 89}


\keyslide{../bilibili_video/slides_p4/slide_0090.jpg}{Slide 90}


\keyslide{../bilibili_video/slides_p4/slide_0091.jpg}{Slide 91}


\keyslide{../bilibili_video/slides_p4/slide_0092.jpg}{Slide 92}


\keyslide{../bilibili_video/slides_p4/slide_0093.jpg}{Slide 93}


\keyslide{../bilibili_video/slides_p4/slide_0094.jpg}{Slide 94}

不是什么别的东西。它并不会真的去驱动工具。那怎么让这串文字真的去驱用一些工具呢？那你需要自己先写好一个小程式。这个小程式是看到语言模型输出要使用工具的指令的时候。

把这串指令真的拿去执行。那等一下我们其实有Colab示范说。这整件事到底是什么意思。我这边先用投影片来说明一下。如果你觉得有点抽象的话。

等一下有Colab实际的说明。告诉你说这一切到底是怎么发生的。好语言模型输出这段文字之后。把这里面的指令取出来。真的执行它。


\keyslide{../bilibili_video/slides_p4/slide_0095.jpg}{Slide 95}


\keyslide{../bilibili_video/slides_p4/slide_0096.jpg}{Slide 96}

真的去呼叫temperature这个函式。真的得到temperature的回应。说摄氏呢32度。那把答案呢放在output里面。那语言模型的使用工具的过程。


\keyslide{../bilibili_video/slides_p4/slide_0097.jpg}{Slide 97}


\keyslide{../bilibili_video/slides_p4/slide_0098.jpg}{Slide 98}

跟得到的答案。其实没必要给使用者看。所以你可以把语言模型这些输出隐藏起来。语言模型实际上还是产生了这些文字。但你没有必要在你的使用者界面上。

告诉使用者你做了这些事。好那语言模型再根据这边的输出。再去做文字接龙。那它就可能可以告诉使用者说。某年某月某日高雄的气温。


\keyslide{../bilibili_video/slides_p4/slide_0099.jpg}{Slide 99}


\keyslide{../bilibili_video/slides_p4/slide_0100.jpg}{Slide 100}

是摄氏32度。使用者还以为这个语言模型做文字接龙。真的能接出正确的温度。但其实不是，背后已经呼叫了工具。

这个语言模型之所以能够回答正确的温度。这是工具的输出。告诉语言模型的。好我们这边呢？实际上来操作一下。

让你真实的体验一下使用工具是怎么回事。那我刚才已经执行了这个colab 在这个colab里面呢？首先我们要先连到huggingface 这个跟上一堂课colab做的事情是一样的。

这边就不再详加说明。然后呢？我们今天使用上周示范程式的最后一部分 pipeline来使用模型。其实上周呢？

讲了好长好长。告诉你用模型有各式各样的方法。什么把模型自己讲的话塞到他嘴巴里等等之类的。其实你在真正用模型的时候 pipeline才是最常用的方式。

就假设你要用huggingface上面的模型的话。其实最容易使用的是pipeline 但之所以上周要讲那么多。是为了让你清楚的了解语言模型背后运作的原理。好那这边我们就使用pipeline呢？

另外呢去下载了gamer3-4b这个模型。好然后我们下载gamer3-4b这个模型以后。接下来呢我们要来让他使用工具。在他使用工具之前。我们先帮他创建两个小工具。

这两个小工具都没什么大不了的。一个工具呢是做乘法。就把multiply 夸号后面给他ab这样的数字。他就回传hab。

有另外一个小工具呢？是做除法给他ab这两个字。有没有ab这样的数字。他就把a除以b 然后回传给你。

好这边有两个工具。但是要注意一下。我们接下来就是要让语言模型呢？去呼叫这两个工具。但不要忘了。

语言模型永远都只能够输出文字。所以他可以输出一段文字叫做multiply 夸号3.4 代表说他好想要用multiply这个工具。把3呢乘上4。

但这就是一串文字。他被输出出来以后还是一串文字。他不会发生任何效果。那在collab上呢？如果你要让一串文字被执行。

你要呼叫一个叫做valuate 叫eval的函式。那你把eval的函式。他后面的夸号里面给他一串文字。他就把那串文字当作指令。


\keyslide{../bilibili_video/slides_p4/slide_0101.jpg}{Slide 101}

真的去执行那串文字里面要模型做的事情。好那所以你用eval加上工具指令。才能真的使用工具。我们现在打eval夸号multiply3.4 那本来multiply3.4这只是一串文字。

但是当我们呢？只用eval夸号multiply3.4的时候。你就真的执行这个工具。就等到这个工具回传的输出。也就是12。

好那接下来就是让语言模型来使用工具吧？那这边使用语言模型的方法呢？这个就是我们上周讲过的。在看语言模型放给他什么样的文字之前。先来看一下使用的套路。

首先你有一个东西叫做message 那message里面就是我们要给语言模型的 prompt 那这边包含了system prompt 这边包含了user的 prompt 那你会发现这边的格式。

跟上周给lama的格式有点不一样。对每个模型他输入的message的格式都是不一样的。你在hacking face上可以找到相关的资讯。所以你要用个模型之前。在hacking face上读一下这个模型的说明。

每个模型他的message的格式可能会略有差异。好把message里面包含了system prompt跟user prompt user prompt丢给pipeline 他就会去执行这个语言模型。他的输出叫做output。

那我们会把output里面的内容。跟语言模型回答有关系的。语言模型输出的部分。语言模型输出的部分呢？把它存在response里面。

然后我们最后再把response印出。就可以看到语言模型输出什么了。那接下来我们来看一下system prompt 跟user prompt里面分别存的什么 user prompt里面我们就是告诉语言模型。

要怎么用一个工具。那这边怎么让语言模型用一个工具呢？我们这边就直接用人看得懂的方法跟它说明。我刚才说每一个工具都是一个函式。怎么用工具呢？

把工具放在tor跟斜线tor里面。那工具的输出呢？会放在toroutput跟斜线toroutput里面。然后呢可用的工具有两个。一个叫multiply。

他会把a乘以b 一个叫divine 他会把a乘以b 然后我的user prompt就是11乘222除以777 正确答案我算过了。

是31.71左右。好，我们真的来执行一下。看看模型给我们什么？看它的输出是什么？

好，他考虑一连串的输出哦。他说他要执行multiply这个工具。得到multiply这个工具的输出。叫24642。

接下来呢？他再根据64242要执行divine这个工具。把24642除以777得到32.258 嗯这个数字怎么怪怪的。我们来验算一下吧？

到底24642除以777是多少呢？哎呀是31.714啊。怎么答案会不一样。难道device这个函式有写错吗？你想想看。

为什么会这个样子。我们想想看哦。到此为止，模型呼叫了工具吗？其实他没有呼叫工具。

他只呼叫了一个寂寞而已。知道吗？他其实并没有操控任何工具。为什么？这一连串的文字。

什么使用工具。使用工具的输出。都是模型自己用文字接龙。输出出来的，所以你今天跟他说你可以用工具。


\keyslide{../bilibili_video/slides_p4/slide_0102.jpg}{Slide 102}


\keyslide{../bilibili_video/slides_p4/slide_0103.jpg}{Slide 103}


\keyslide{../bilibili_video/slides_p4/slide_0104.jpg}{Slide 104}


\keyslide{../bilibili_video/slides_p4/slide_0105.jpg}{Slide 105}


\keyslide{../bilibili_video/slides_p4/slide_0106.jpg}{Slide 106}

好，那他就使用了工具。他就输出一段他想要使用工具的请求。但是这段请求并没有真的变成执行工具的指令。他就是一串文字而已。


\keyslide{../bilibili_video/slides_p4/slide_0107.jpg}{Slide 107}

\textbf{幻灯片引用链接:}
\begin{itemize}
  \item \href{https://www.anthropic.com/news/3-5-models-and-computer-use}{https://www.anthropic.com/news/3-5-models-and-computer-use}
\end{itemize}
\begin{itemize}
  \item \href{https://openai.com/zh-Hant/index/introducing-chatgpt-agent/}{https://openai.com/zh-Hant/index/introducing-chatgpt-agent/}
\end{itemize}

那语言模型呢看到这段文字。他在自己继续去做文字接龙。那既然有人请求了工具。现在发生什么事呢？当然是要得到工具的输出。

然后Gamer实在是非常厉害 11×222 他在没有工具的前提下。他居然心算算对了呀。是24642。

实在是非常厉害。好，那有了这个相成的结果以后呢？再把24642除以777 看看是多少？

那这边仍然没有呼叫任何工具。一切都是文字接龙产生的结果。所以根据24642除777 再去文字接龙。接出个32.258。

这是语言模型心算的结果。你要惊叹的事情是。他居然跟正确答案只差了1而已。只差了一点多而已。非常的厉害。

但到目前为止。语言模型，他并没有呼叫任何工具。所以在上面那段程式码中。语言模型只用了寂寞。


\keyslide{../bilibili_video/slides_p4/slide_0108.jpg}{Slide 108}

\textbf{幻灯片引用链接:}
\begin{itemize}
  \item \href{https://www.anthropic.com/news/3-5-models-and-computer-use}{https://www.anthropic.com/news/3-5-models-and-computer-use}
\end{itemize}
\begin{itemize}
  \item \href{https://openai.com/zh-Hant/index/introducing-chatgpt-agent/}{https://openai.com/zh-Hant/index/introducing-chatgpt-agent/}
\end{itemize}

他没有用任何的工具。怎么让语言模型真正的去呼叫工具呢？所以以下是让语言模型呼叫工具真正的方法。我们先来看一下这张图的说明。好我们给予语言模型 system prompt。


\keyslide{../bilibili_video/slides_p4/slide_0109.jpg}{Slide 109}

我们给他 user prompt 接下来他要做文字接龙。好他要说什么呢？他可能会根据这些输入的结果。根据这些 context 的结果。

觉得他应该要用个工具。他就说好那我要用工具。但用完工具以后他其实会继续接龙下去。他会接这些工具的输出。这是一个幻觉。

这是一个 hallucinate 的结果。不要管他把他说要用工具的请求结出来。也就是把Tor跟Tor斜线中间工具的指令拿出来。丢给EVAL这个函式。我们有说工具的指令是一串文字啊。

他没办法真的被执行。丢给EVAL这个函式才能够真的在collab上执行这个函式。得到工具的输出。那至于Tor后面 hallucinate 出一大堆。他想像用了这个工具以后。

应该会有多美好的那些字句。不要理他就把他丢掉。好所以现在的状况是。执行语言模型呼叫了一个工具。得到了工具的输出。

但之所以得到工具输出。是我们帮他去真的执行了这个工具。好那把工具的指令跟工具的输出。都放到语言模型的context 所以现在的对话变成。

人问了一个问题。语言模型呼叫一个工具。语言模型得到一个工具的输出。那有一些模型会有一些特殊的栏位。专门放工具的输出。

不过Gemma没有这个栏位。他要么是systempump 要么是。其实啊，告诉你一个秘密啊。

Gemma没有真正的systempump 他在输入的时候你可以给他systempump 但他还把systempump 贴在第一个userpump前面这样子。不知道就是一个不知所欲的动作。

很神奇，他其实没有真正的systempump 好。那他只有systempump userpump。

跟assistant讲的话。那他没有说工具讲的话。所以我只好把工具讲的话。当作使用者讲的话。然后我发现说呢？

如果你把工具的输出。当作使用者讲的话。很多时候呢？这个Gemma会误判。以为他是真正的使用者讲的话。

所以我这边要特别强调。它是tor output跟斜线tor output 他读到这两个符号。他会知道说，这个是一个工具的输出。

让语言模型呢？再根据上面的这些context 继续去做文字揭楼。那他想要用第二个工具。产生第二个工具指令。

然后呢？再把第二个工具指令。并丢给EVAL 产生工具的输出。那工具指令后面。

不管他要 hallucinate 出什么东西来。都不要理他，接下来，把第二次用工具的指令。

跟第二次工具的输出。再放到对话历史记录。所以现在对话变成。使用者的要求。第一个工具指令。

第一个工具输出。第二次用工具。第二次工具输出。再继续去做揭楼。当他最后揭出来的结果。

没有使用工具的要求的时候。如果使用工具的要求。就会出现拓跟斜线拓。没有出现拓跟斜线拓。就代表说他觉得不需要用工具。

那这个就是语言模型最终的答案。那中间这些呼叫工具。还有工具的输出。你其实没必要给使用者看。使用者看到语言模型最终输出的答案。

就可以了，好，那个真的来执行上面那段程式码吧？好，所以这边做的事情是这样子的。我们给语言模型使用工具的quamp 我们给语言模型的使用者的输入。

然后把它放到message 根据这个GEMMA固定的格式。存到message里面。然后接下来呢？我们会进入一个坏油回圈。

在这个坏油回圈里面。我们每次开始的时候呢？都会把这个message丢给pipeline 让它得到一个输出。就是我们会跑语言模型去做文字接龙。

把它揭出来的结果放在response里面。那如果response里面有要用工具的指令。我们就把tour跟斜线tour里面的指令结出来。其他地方就不要管它。其他地方是alucination。

把tour跟tour 斜线tour之间的指令结出来。然后我们会把结出来的。互交工具的指令印出来给你看。然后呢？

接下来再把这个指令。我们存在comment里面。把这个comment用EVAL执行。然后它执行出来。结果把它转成那个字串。

然后把它存在touroutput这个变数里面。然后我们会把工具的输出。也印出来给你看。不过这样的步骤。在真的做一个平台的时候。

你可以隐藏起来。不让使用者看到。那接下来呢？我们就把这个刚才呢？模型呢？

使用工具这件事情。放到它历史的对话里面。它才知道它刚才用了工具。我们再把工具执行的结果。也放到message里面。

放到历史的对话里面。模型才知道它用了工具。又得到了工具的结果。然后这个回圈就继续。它就会再继续做接龙。

看看它有没有要用工具。如果有要用工具。那就会反复刚才的步骤。帮模型执行工具。如果它没有要用工具。

如果输出的结果里面。没有跟工具有关的符号。那llm就会真的输出它的答案。我们就把llm文字接龙的结果。当作最终的答案。

印出来给使用者看。好，那我们来看看发生什么事吧？呼叫工具 multiply。

然后得到工具的回传。然后再呼叫一次工具。它呼叫divine 因为我们要把24642除以777 得到31.714。

然后最后模型输出的结果就是 111乘22除以777等于31.714 这一串。是最后使用者看到的内容。前面呼叫工具。

等等不一定要给使用者看。以前那个ChangeVision的平台。它呼叫工具的时候。它都会告诉你说我要做什么？我要做什么？

现在它往往都不告诉你要做什么？它会把用工具的行为传起来。所以我现在有时候都怀疑说。它到底是不是这个答案。到底是它自己产生的。

还是偷偷调用某个工具所产生的。所以对使用者来说。它看到的就是这一串输出。它就会觉得这个模型很厉害。会做加减乘除。

用文字接容做加减乘除。居然还没做错。其实这是呼叫了一个工具以后的结果。好我们再做另外一个工具。这个工具是会给我们温度。

这个工具做的事情是这样子的。你就给它一个城市的名字。跟一个时间，然后它就会说。这个城市在某个时间是摄氏30度。

那当然这是个假的工具。如果你真的要用有用的工具的话。也许你可以去连一下气象局的API 让它真的传一个真的在某个时间的气温给你。那我们这边是做一个假的工具。

好在这个假的工具执行以后呢？它就会跟你说高雄的气温是摄氏30度。好那我们现在就让模型真的来用一下。这个工具试试看吧？然后我们现在问的问题改成。

告诉我高雄1月11号的天气如何？然后看看模型呢会有什么样的反应。好然后呢他就呼叫工具。呼叫GetTemperature这个工具。他知道城市呢是高雄。


\keyslide{../bilibili_video/slides_p4/slide_0110.jpg}{Slide 110}


\keyslide{../bilibili_video/slides_p4/slide_0111.jpg}{Slide 111}

时间是1月11号。工具告诉他现在气温是30度。然后所以他真正最后使用者看到的输出就是。好的高雄1月11号气温是30度。使用者还以为模型中文字接容真的能够接出气温来呢？

其实不是，这是工具告诉模型的。啊其实今天这些语言模型都蛮厉害的哦。你告诉他，如果你给他怪怪的东西。

告诉他现在呢？温度非常非常高。这个一定比太阳表面的温度还要高。我们再让语言模型使用这个工具看看。他会说些什么？


\keyslide{../bilibili_video/slides_p4/slide_0112.jpg}{Slide 112}

来看他今天说哇高雄气温居然这么高。这有点离谱，你是不是输入有误了。所以很厉害哦。他并不是工具输入的结果。


\keyslide{../bilibili_video/slides_p4/slide_0113.jpg}{Slide 113}

他就一定会相信如果工具输出荒唐的结果。他是有机会知道工具的输出有问题的。当然语言模型其实只在你需要用工具的时候才用工具啦。像我们刚才都问数学问题或是问气温那需要工具。但假设我们现在就是问他比如说你好吗？


\keyslide{../bilibili_video/slides_p4/slide_0114.jpg}{Slide 114}

看看他在说什么？所以语言模型可以自己决定他要不要用工具。我问他你好吗？那就不会呼叫工具。因为没有工具是做这件事的。

他说我很好谢谢这样。所以这个就是语言模型使用工具的范例。那我们刚才讲了语言模型呢？可以怎么使用工具。那在语言模型可以使用的工具里面。

有一个最通用也最强大的。其实就是使用电脑。我们很多人把它叫做computer use 由现在呢？语言模型可以直接用滑鼠跟键盘。

去操控一台电脑。他可以给他荧幕的画面。给他任务的指示。按照荧幕画面跟任务的指示。他就去操控你的滑鼠跟键盘。

那他可以操控滑鼠跟键盘。有什么厉害的地方呢？他基本上就跟一个人类没什么差别了。人类能用电脑做的是。现在语言模型也有机会可以直接用电脑来做。

那讲到这个computer use 这个Cloud呢？跟ChatGPT呢？都出了computer use的功能。那我这边呢可以给大家看一下。

当ChatGPT在用电脑的时候。看起来像是什么样子。那我们直接呢？把ChatGPT开起来给大家。好我今天只是想给大家看一下说。

现在ChatGPT有一个代理人模式。那代理人模式呢？翻译成英文就是agent mode 那等一下会再更详细讲一下 agent是什么意思。

那当你使用代理人模式的时候。就有一定的可能性。他会开启一个萤幕画面开始做事。比如说我这边跟语言模型说。帮我订高铁票。

9月20号上午9点到10点。台北到左营两张票。那他要怎么订票呢？他会开启一个萤幕画面。然后呢就开始订票。

我们从头播放一下。让你知道说语言模型在做什么事。好那语言模型呢？如果是ChatGPT 他并不会真的去操控你的。

电脑，他会做的事情是。他开启了一个萤幕画面。所以你可以想像说。有一个电脑呢？

在云端ChatGPT 就用那一台电脑。用给你看，所以你看到了。他就现在就搜寻。

台湾高铁相关的资讯。为了要订高铁票。找到高铁网站。好那我们刚才说要做什么呢？我们刚才说。

从台北到左营。所以他就先点。台北，点，点到。

好那接下来呢？要把目的地呢？设为左营，好所以点，目的地。

然后他知道设为左营。他会一边做事。一边讲话，该告诉你说，他现在想要做什么？

然后说要两张票。所以他只要设两张票。点地一下，没点到，好。

往下移，滑鼠往下移一点。点到了，然后看到一些数字。他决定选2。

选到2了，好接下来呢？再设定时间，时间呢？是9月20号。

好那看到日历。他点了一下，要设20号，但不知道为什么？点到19。

看了一下萤幕画面。发现是19 发现这个任务没有完成。所以再点一次。还是没点到。

再点一次，点到了，这次记得，要设20号，能够成功设到20号吗？

还是没有，然后再一次，再点开，注意哦，要设到20号。

设到了，接下来呢？要输入验证码。然后他还输入时间。输入一下时间。

这个9点到10点。所以要选个9点。没问题，剩下输入验证码了。他不肯输入验证码。

所以就结束了。这个，你自己订票，一定是快很多了。那你可以想象说。

以后如果这个功能。真的完善的话。不会像刚才，像点半天，点不到的话。

这个实在是威力无穷。语言模型，是怎么操控，滑鼠跟键盘的呢？其实。

就跟刚才使用工具。是一样的，滑鼠跟键盘，也就是工具，所以。

当你要叫模型。去使用电脑的时候。你真正做的事情就是。给他一个萤幕画面。然后跟他讲说。


\keyslide{../bilibili_video/slides_p4/slide_0115.jpg}{Slide 115}


\keyslide{../bilibili_video/slides_p4/slide_0116.jpg}{Slide 116}


\keyslide{../bilibili_video/slides_p4/slide_0117.jpg}{Slide 117}

你现在可以用的工具是。有键盘，你可以输入你要的字幕。你可以移动滑鼠。到萤幕上的某一个位置。

你可以点击滑鼠。左键到右键，然后叫他决定。他接下来，要做什么？

真的能透过这样子。让语言模型使用电脑吗？我这边就真的示范一下。那开启另外一个呢？我跟ChairGPT的历史对话记录。

这边呢？我就是给他一张图片。给他一个萤幕截图。然后告诉他要怎么使用键盘跟滑鼠。然后问他。


\keyslide{../bilibili_video/slides_p4/slide_0118.jpg}{Slide 118}


\keyslide{../bilibili_video/slides_p4/slide_0119.jpg}{Slide 119}

如果你要订阅李鸿毅老师的YouTube频道。那下一步是什么？哇他知道滑鼠移到640 340这个位置。那我真的实际。我真的实际试了一下。


\keyslide{../bilibili_video/slides_p4/slide_0120.jpg}{Slide 120}

640到340 真的就是大概在这个搜寻栏的位置。所以他知道搜寻栏在哪里？好再来问他说。那下一步呢？

他要点滑鼠的左键。在下一步呢？键盘按李鸿毅YouTube 在下一步呢？按Enter。

所以他完全知道要怎么搜寻一个YouTube频道。所以这个Computer Use的原理大概就是这个样子。不过不是每个模型都能够真的做到啦。我刚才要开GPT5的Thinking Mode 才有办法答对。

那如果是用更旧的模型。比如4.1 他就没有办法真的答对刚才的指令。好那刚才讲了好多东西。都是人类外部强放进Context里面的。


\keyslide{../bilibili_video/slides_p4/slide_0121.jpg}{Slide 121}

现在模型啊，它的Context里面。可能还可以包含。自己产生的思考过程。现在有很多模型。

号称可以做深度思考。比如说ChagVT的O系列模型啊 D-Sig的R系列模型啊 Germanline的D-Think啊。都是号称说模型可以做深度思考。

所以深度思考的意思是。当你问一个模型的时候。他先不直接给你答案。而是他会先演一个脑内小剧场。所以他会演一个脑内小剧场。

那如果是个数学问题的话。往往他就会把各种不同的解法。都在脑中演练一遍。他说先来解个。先用A解法解看看。

验算一下答案。嗯不对，试一下B解法。好像是对的，那C解法也可以试一下。

嗯看起来不太好。他会演一个脑内小剧场。把各种不同的可能性。在脑中考虑一次给你看。然后考虑完之后。

他再根据刚才考虑的结果。来真的解题，给你一个最终的答案。那在脑内小剧场中。这个模型可以做的事情。

包括比如说，规划要怎么解题。他可能会做各种不同的尝试。他也会验证每次尝试的答案。那今天啊。

这个脑内小剧场对使用者来说。你是可以选择不看的。甚至你根本就不能看。像ChairGVTR 他基本上是不让你看。

这个脑内小剧场的。他只是给你看。脑内小剧场的摘要。所以你并不知道。他完整的脑内小剧场里面。


\keyslide{../bilibili_video/slides_p4/slide_0122.jpg}{Slide 122}


\keyslide{../bilibili_video/slides_p4/slide_0123.jpg}{Slide 123}

到底演了什么样的东西。所以这一些模型。自己产生的思考过程。其实也可以看作是 context的一部分。

只是context 不是我们人外加给他的。而是他自己放进去的。这个模型，就是根据他自己产生的context。

再来做文字接龙。再产生最终的答案。那如果你想知道。模型是怎么学会做深度思考的。那在上个学期的。


\keyslide{../bilibili_video/slides_p4/slide_0124.jpg}{Slide 124}

\textbf{幻灯片引用链接:}
\begin{itemize}
  \item \href{https://youtu.be/bJFtcwLSNxl?si=hT-nKkFEEd-zJypH}{https://youtu.be/bJFtcwLSNxl?si=hT-nKkFEEd-zJypH}
\end{itemize}


\keyslide{../bilibili_video/slides_p4/slide_0125.jpg}{Slide 125}


\keyslide{../bilibili_video/slides_p4/slide_0126.jpg}{Slide 126}

这个机器学习的课程的第七讲。有讲说像D6 R系列这种模型。是怎么练成深度思考能力的。那我把这个影片连结留在这边。

如果你有兴趣的话。给大家参考，好，所以我们讲了一连串 context里面要包含什么？

我们说要有User Pump 要有System Pump 要有对话历史记录。也要有长期的记忆。要有一些搜寻引擎给我们的结果。


\keyslide{../bilibili_video/slides_p4/slide_0127.jpg}{Slide 127}

也要有工具使用的结果。也要有Reasoning的结果。它非常非常的长。所以Context Engineering的核心目标是什么？它核心目标就是一句话。

避免塞爆Context 想办法只放需要的东西进入Context 清理掉不需要的内容。那在AI Agent的时代 Context Engineering尤其非常重要。

那我们先来介绍一下AI Agent 让你更能够体会说。为什么在AI Agent的时代 Context Engineering 会是一个关键的技术。

那过去我们一般使用AI的方式。就是一问一答。你问他一个问题。语言模型接融出一个答案。那你得到答案以后你就满意了。


\keyslide{../bilibili_video/slides_p4/slide_0128.jpg}{Slide 128}

那今天很多时候。我们会更进一步采取Agentic Word 也就是说帮语言模型定一个SOP 让它按照SOP来执行。那这特别适合一些比较复杂。

但你知道有哪些步骤的任务。比如说如果我们让语言模型批改作业。那直接批改作业。直接给语言模型的作业。问他说这个作业应该拿几分。

那往往会有各式各样的问题。所以批改作业的流程。你可能要设计成多个步骤。作业输入以后。先要检查这个作业。

到底是不是真正的作业。还是有人呢？想要对这个，模型做这个 Pump Injection Attack。

所以第一步先检查。是不是真正的作业。第二步再给分数。那最后可能还要验证说。这个分数给的。

是不是合理，所以如果今天一个任务有 SOP 那你可以让语言模型分多个步骤来完成。那这样可以让语言模型完成更复杂的任务。


\keyslide{../bilibili_video/slides_p4/slide_0129.jpg}{Slide 129}


\keyslide{../bilibili_video/slides_p4/slide_0131.jpg}{Slide 131}


\keyslide{../bilibili_video/slides_p4/slide_0132.jpg}{Slide 132}

\textbf{幻灯片引用链接:}
\begin{itemize}
  \item \href{https://github.com/go}{https://github.com/go}
\end{itemize}

AI Agent又是更进一步。让语言模型自己决定解决问题的步骤。而且根据在解题过程中发生的变化。语言模型应该要能够灵活调整它的计划。比如说你叫语言模型解决某一个研究问题。


\keyslide{../bilibili_video/slides_p4/slide_0133.jpg}{Slide 133}


\keyslide{../bilibili_video/slides_p4/slide_0134.jpg}{Slide 134}


\keyslide{../bilibili_video/slides_p4/slide_0135.jpg}{Slide 135}

它可能先上网搜集资料。得到资料之后。它可能可以形成自己的假设。有了假设以后就要做实验来验证。虽然这边放了一些烧屏之类的。

不过如果是Computer Science的话。做实验通常指的是写程式。那对语言模型来说。写程式完全不成问题。那写完成程式以后呢？

那语言模型呢？再根据程式验证的结果。看看假设有没有被验证。如果不成立，再上网搜寻资料。

重新寻找假说。曾经形成假设。所以这可能是一个。没有办法事先预期要有哪些步骤的问题。语言模型可能要进行多个循环。

一边收集文献。一边做实验，那可能要执行很多不同的步骤。最后才能够写出一个技术报告。或写出一本论文给人类。

所以AI Agent要模型做的事情。是自己决定解题的步骤。并且在解题的过程中。在执行任务的过程中。灵活的根据现实的情况。

来调整它的规划。那AI Agent整体看法。可以看成是一个循环。什么样的循环呢？首先人类先给语言模型一个目标。

然后语言模型根据现在的状况。现在的输入，现在的状况，现在的输入，可以叫做observation。


\keyslide{../bilibili_video/slides_p4/slide_0136.jpg}{Slide 136}

它根据observation 去采取一个行为。比如这个行为可能是调用工具。这个行为可能是写一段程式码。这个行为可能是输出一个报告。

这个行为可能是跟人类联络等等。好的，语言模型输出一个行为之后。那这个行为可能是跟某个工具互动。可能是跟人类的使用者互动。

那人类啊工具啊这些语言模型以外的东西。统称为环境，我们这边用一个地球来表示环境。那环境呢？收到语言模型的要执行的行为之后。

环境呢就会有所改变。可能给一个对应的输出。比如说语言模型。请求人类提供多一点资讯。人类提供多一点资讯。

语言模型呢执行了某一个工具。工具给一个输出。总之语言模型看到的observation 就改变了。有了新的observation。

语言模型就会去采取新的行为。新的行为有了新的行为。就会又有不一样的observation 以此类推。以此就不断的循环下去。

好那AI agent 相较于传统的agent 有什么样的优势呢？比如说AlphaGo 可以看作是一个传统的agent。

他可以下围棋。对他来说，他的环境就是他的围棋的对手。他的目标就是要引起。但是用语言模型来运作一个AI agent。

跟传统AI agent不一样的地方是。语言模型它输出的是文字。所以你可以想成它的输出。有近乎无限的可能 AlphaGo可以有的输出是事前设定好的。

它就是只能够在棋盘上。选择一个落子的位置。但语言模型有近乎无限输出的可能。然后语言模型呢？可以用人类的语言提供回馈。

你可以用人类语言。要求这个模型采取某些行为 AlphaGo听不懂人话。如果你叫他要执行某一个行为。落子在某一个地方。

你其实做不到的。但是现在的语言模型可以。你可以用文字来控制它。可以用文字来提供回馈。那刚才我们其实有看到。

Gemini呢有一个agent mode 不过我知道Gemini的agent mode 要付费才能够使用。所以不是每个人都可以体验。今天的AI agent长什么样子。

那如果你想玩一个免费的 AI agent的话。那我觉得你可以考虑 Gemini CLI 它是免费的。

然后呢？用Gemini CLI 你背后还可以免费。偷耗Gemini 2.5 Pro 它背后呢？

有时候会用Gemini 2.5 Pro运作。然后它是免费的。所以等于是可以偷偷用 Gemini 2.5 Pro 我们可以实际跑一下。

Gemini CLI给大家感受一下。看看它是什么样子。好那至于安装。大家就在自己研究。你安装好以后呢？


\keyslide{../bilibili_video/slides_p4/slide_0137.jpg}{Slide 137}


\keyslide{../bilibili_video/slides_p4/slide_0138.jpg}{Slide 138}


\keyslide{../bilibili_video/slides_p4/slide_0139.jpg}{Slide 139}

在这个指定的界面上打Gemini 然后等于是把Gemini呼叫出来。希望我可以成功的呼叫它。好成功的呼叫出来了。好接下来呼叫出来以后。

你要它做什么都可以。然后告诉你 Gemini CLI非常危险。为什么它非常危险呢？因为它是真的会碰触你的电脑。

跟刚才看到的 ChairGVT的那个Agent Mode不一样 Agent Mode是用一个 Agent Mode是用一个线上的电脑。所以它不会影响你的电脑。

所以它不会影响你的电脑 Gemini会真的碰触你的电脑。比如说，我们可以叫它。关闭开启的投影片。

我们试试看喔。看看它关不关喔。好，它开始执行喔。它会不会做这件事呢？

它有时候做有时候不做啦。它还好在做之前呢？它还会问你，试着关起来咯。哇。

投影片被关起来了。然后，所以如果我等下叫它关闭直播。我们直播就关闭了。这真的很危险啦。

不过我试过叫它关电脑。它是不关的，不过你叫它删档案。它是愿意删的。所以。

哇，这个Gemini CLI用起来。真的是很危险。好，那你要叫它做什么？

就叫它做，任何事，比如说，我们叫它，做一个贪食蛇的游戏。

然后呢？这不是一个，一个步骤用Pump 就可以完成的事情。你就会看到它。

是一个agent 它会花多个步骤完成。然后右下角，要告诉你说，现在context用了多少？

Gemini的context 真的非常长喔。而且它采用了1\%的context 还有99\%的context 你还没把它的context。

占满，好，跟它说，我们要做一个贪食蛇的游戏。然后你看这边。

它就告诉你它在做什么？它在develop game的logic 然后呢？好，它说它有一个计划。

那这些模型，通常做计划之前呢？它需要人同意啦。他们还是比较小心的。所以说同意。


\keyslide{../bilibili_video/slides_p4/slide_0140.jpg}{Slide 140}


\keyslide{../bilibili_video/slides_p4/slide_0141.jpg}{Slide 141}


\keyslide{../bilibili_video/slides_p4/slide_0142.jpg}{Slide 142}

好，开始来写这个程式喔。然后这边就，它要建立一个资料夹。那你这边可以让它allow always。


\keyslide{../bilibili_video/slides_p4/slide_0143.jpg}{Slide 143}

之后就不会再问你了。不过我们谨慎起见呢？还是人来按个enter 做一个资料夹。接下来产生三个档案。

index.html sell.css script.js 那反正它自己会开启档案。自己产生输出。

自己把档案存起来。这它自己都会做。我们就交给Gemini来完成。那这边我们就一直按allow 都可以。

它就做什么都可以。它就做了，它就开始写程式。做完了吗？做完了。

能不能玩呢？等一下，第一个问题就是。这个游戏到底存在哪里？它应该是在这个资料夹下执行的。

然后呢？刚才不是make一个snake game 在这里。应该是点这个index就可以玩了。按enter开始游戏。


\keyslide{../bilibili_video/slides_p4/slide_0145.jpg}{Slide 145}

我是绿色的蛇啊。我还以为我是红色的东西呢？这还让不让人好好玩呢？可以可以可以。可以它太快了。

还让不让人好好玩。就这样，好，那这个太快了。我觉得。

这没办法好好玩。我跟他说，太快了，蛇太快了，没办法好好玩。

好，他就改了，他知道我这句话的意思。他就改了，他知道我这句话的意思。

他就改了，他知道我这句话的意思。他就改了，他知道我这句话的意思。他就改了。

看，他把100改成150 哦。这边都改了，哦。


\keyslide{../bilibili_video/slides_p4/slide_0146.jpg}{Slide 146}


\keyslide{../bilibili_video/slides_p4/slide_0147.jpg}{Slide 147}


\keyslide{../bilibili_video/slides_p4/slide_0148.jpg}{Slide 148}


\keyslide{../bilibili_video/slides_p4/slide_0149.jpg}{Slide 149}


\keyslide{../bilibili_video/slides_p4/slide_0150.jpg}{Slide 150}


\keyslide{../bilibili_video/slides_p4/slide_0151.jpg}{Slide 151}

好，他说改好了，我们再来玩一次吧？希望蛇真的有变慢。按enter开始游戏。

哦，感觉真的有变慢了。哎，这太难了，哎呀。

这没办法好好玩。不过就是，真的有变慢了。这太难了，哎呀。

这没办法好好玩。不过就是，这个感觉，你现在了解说。以后写程式。

就是这个样子了。这个就是vibe coding 好。那展示完以后。我们把投影片开回来吧？


\keyslide{../bilibili_video/slides_p4/slide_0152.jpg}{Slide 152}

到底能不能够用 Gemini 成功把投影片开回来呢？不一定能成功。把刚才关闭的投影片再打开。看看他做不做。

他不一定能够知道刚才关起来的是什么？我来看一下哦。他知不知道刚才关了什么？然后能不能再打开。嗯。


\keyslide{../bilibili_video/slides_p4/slide_0153.jpg}{Slide 153}


\keyslide{../bilibili_video/slides_p4/slide_0154.jpg}{Slide 154}

他搜寻我的投影片。他看看有哪些投影片。他问我打开哪个？所以他显然不知道。没关系。

我告诉他 NL资料夹下。那个 Agent B6。


\keyslide{../bilibili_video/slides_p4/slide_0155.jpg}{Slide 155}


\keyslide{../bilibili_video/slides_p4/slide_0156.jpg}{Slide 156}

\textbf{幻灯片引用链接:}
\begin{itemize}
  \item \href{https://github.com/go}{https://github.com/go}
\end{itemize}

好，看看他知不知道我的意思哦。哎，感觉应该可以。好又可以上课了。

好，又可以上课了。所以，就是这么神奇啊。好。

那刚才呢？展示了 Gemini CLI的运作。所以了解说，啊今天这个 AI Agent。

有多么强大，好，那我们从，语言模型的角度来看 AI Agent。


\keyslide{../bilibili_video/slides_p4/slide_0157.jpg}{Slide 157}


\keyslide{../bilibili_video/slides_p4/slide_0158.jpg}{Slide 158}


\keyslide{../bilibili_video/slides_p4/slide_0159.jpg}{Slide 159}

他到底在解什么样的问题。从语言模型的角度来看。其实他做的事情。始终都是文字接龙。一开始有一个System Point。

那里面可能写了。人要他做的事情。然后呢？他会有第一个Observation 看看现在四周发生了什么事情。

那他采取了一个合适的回应。采取了一个Action 然后呢？这个Action 会影响环境。

所以会看到，不一样的Observation 语言模型。再根据新的Observation 再采取下一个Action。

下一个Action 又会影响环境。这个步骤就反复执行下去。但是语言模型。实际上做的事情。

始终都只是文字接龙。他根据Observation 1 产生出Action 1 根据Observation 1 Action 1跟Observation 2。

产生Action 2 以此类推。他做的事情，始终都是文字接龙。跟在做对话。

其实是没有什么本质上的区别的。所以其实 AI Agent 他其实不太像是一个新的技术。他比较像是。


\keyslide{../bilibili_video/slides_p4/slide_0160.jpg}{Slide 160}

依靠语言模型。它就是很会接龙。所以他可以做长时间的运行。他可以帮你做更复杂的任务。那运行AI Agent的挑战是什么呢？

他最大的挑战之一就是。我们今天要语言模型做的事情。可能非常的复杂。可能有非常多的步骤。如果你要语言模型。

帮你写一本博士论文。那显然他需要做非常多的实验。才有办法最终写出一本博士论文。当输入过长的时候。语言模型可能会发疯。

他就没有办法好好的做文字接龙。那今天各个你可以用的模型。往往都会告诉你说。这个模型的context window size有多大 context window size就是。

这个模型可以输入的长度上限。你可以想成说。这些模型在训练的时候。他就只能看最长这么长的输入做接龙。再更长会发生什么行为。


\keyslide{../bilibili_video/slides_p4/slide_0161.jpg}{Slide 161}

\textbf{幻灯片引用链接:}
\begin{itemize}
  \item \href{https://www.meibel.ai/post/understanding-the-impact-of-increasing-li-cc}{https://www.meibel.ai/post/understanding-the-impact-of-increasing-li-cc}
\end{itemize}


\keyslide{../bilibili_video/slides_p4/slide_0162.jpg}{Slide 162}

\textbf{幻灯片引用链接:}
\begin{itemize}
  \item \href{https://www.artfish.ai/p/long-context-Ilms}{https://www.artfish.ai/p/long-context-Ilms}
\end{itemize}

没有人知道，他可能会开始发疯。接触各式各样奇奇怪怪的东西。你不知道会发生什么样的事情。那今天各个知名的语言模型。

它的context window 也就它可以当做输入的这个context 可以放多少的token呢？这边每一个点代表一个语言模型。横轴代表的是时间。

纵轴代表的是context里面可以放几个token 所以你可以看。当年GPT4 可以放3万多个token 蛮长的。

后来 Cloud系列出来。号称可以放10万个token 大家都惊呆了。可以放10万个token。

后来GEMINI 说它可以放百万个token 大家都惊呆了。居然可以放百万个token 后来LAMA4最近出的。

号称可以放千万个token 为什么会需要放这么长的token呢？因为如果要让语言模型变成A-ZO 长时间运作。它就要看非常长的历史记录。

所以它确实需要非常长的输入。像GEMINI 1.5 它号称可以输入200万个token 200万个token到底是什么样的概念呢？它可以读完哈利波特全集。

加几乎读完三本魔戒。所以这是一个非常长的输入。大家可能会说。那GEMINI或者是其他模型。有这么长的输入。

那不就没事了吗？反正就是让模型凭借它接龙的能力。去长时间运行。它就变成了一个AI agent 但现在的问题是。

能输入上百万个token 并不代表能读懂上百万个token 就好像你把哈利波特全集翻一遍。然后你真的记得全部的内容吗？你可能没办法记得所有的内容。

语言模型也是一样。它输入说支持可以支持200万个token 并不代表200万个token 里面发生的事情。它通通都清清楚楚。

有太多实验告诉你说。我们在还没有到达语言模型。输入长度上限的时候。语言模型就已经开始对于输入感到困惑了。举例来说。

这边是一篇DataBrick他们发的博文。他们说，大家都说要做RIG RIG就是要让语言模型去上网搜一些资料。再把搜寻的资料放到context里面。


\keyslide{../bilibili_video/slides_p4/slide_0163.jpg}{Slide 163}

\textbf{幻灯片引用链接:}
\begin{itemize}
  \item \href{https://www.databricks.com/blog/long-context-}{https://www.databricks.com/blog/long-context-}
\end{itemize}

搜到的资料越多越好吗？直觉上搜到资料应该是越多越好吧？但他们发现说。对于不同的语言模型而言。这边横轴是把多少搜寻到的token。

放到context里面。重轴是语言模型回答的正确率。如果今天输入还没有很长的时候。多数状况下，这边每一条线都是不同的语言模型。


\keyslide{../bilibili_video/slides_p4/slide_0164.jpg}{Slide 164}

\textbf{幻灯片引用链接:}
\begin{itemize}
  \item \href{https://www.databricks.com/blog/long-context-}{https://www.databricks.com/blog/long-context-}
\end{itemize}

多数状况下，资料越多结果越好。这很直觉，搜寻到的资料越多。当然里面越有可能包含正确答案。

语言模型越有可能给我们正确的回答。但是假设资料太多。语言模型就开始发疯了。看不下去啊，就算里面有答案也读不了啊。

没有办法正确的回答问题。所以看很多模型。它的正确率都是先升后降。像绿色的这个。什么meta405b。

给它最多资料的时候。还比不上只给它最少资料的状况的。给它一堆资料。它头晕目眩，根本没有办法回答问题。


\keyslide{../bilibili_video/slides_p4/slide_0165.jpg}{Slide 165}

\textbf{幻灯片引用链接:}
\begin{itemize}
  \item \href{https://arxiv.org/abs/2307.03172}{https://arxiv.org/abs/2307.03172}
\end{itemize}


\keyslide{../bilibili_video/slides_p4/slide_0166.jpg}{Slide 166}

\textbf{幻灯片引用链接:}
\begin{itemize}
  \item \href{https://arxiv.org/pdf/2505.06120v1}{https://arxiv.org/pdf/2505.06120v1}
\end{itemize}


\keyslide{../bilibili_video/slides_p4/slide_0167.jpg}{Slide 167}

\textbf{幻灯片引用链接:}
\begin{itemize}
  \item \href{https://arxiv.org/pdf/2505.06120v1}{https://arxiv.org/pdf/2505.06120v1}
\end{itemize}


\keyslide{../bilibili_video/slides_p4/slide_0168.jpg}{Slide 168}

\textbf{幻灯片引用链接:}
\begin{itemize}
  \item \href{https://arxiv.org/pdf/2505.06120v1}{https://arxiv.org/pdf/2505.06120v1}
\end{itemize}


\keyslide{../bilibili_video/slides_p4/slide_0169.jpg}{Slide 169}

\textbf{幻灯片引用链接:}
\begin{itemize}
  \item \href{https://research.trychroma.com/context-rot}{https://research.trychroma.com/context-rot}
\end{itemize}


\keyslide{../bilibili_video/slides_p4/slide_0170.jpg}{Slide 170}

\textbf{幻灯片引用链接:}
\begin{itemize}
  \item \href{https://research.trychroma.com/context-rot}{https://research.trychroma.com/context-rot}
\end{itemize}

这边有另外一篇paper 这边paper里面观察到一个现象。叫Loss in the Middle 就是语言模型对于输入的context 它通常比较记得开头跟结尾。

这边paper里面做的实验是这个样子的。我们在做RIG的时候。给模型20篇搜寻到的文章。长度大概4000个token左右。那这边是一个比较早的文章。

所以它做GPT 3.5 那4000个token 当时也已经算是蛮长的context 然后这边的横轴代表说。这个正确答案出现在第几篇文章。

他发现说如果正确答案出现在最前面的文章。跟context最后面的文章。模型比较有机会答对问题。出现在中间模型就很容易答错。而且这边这个红色的虚线。

红色的虚线代表说。模型凭借着自己的知识。就是根本就不给它context 不做RIG 凭借自己原来有的知识去做GPT。

如果你今天正确答案是出现在。一大堆文章的中间中间的文章。那还比不上让模型直接回答呢？模型直接回答正确率还比较高。所以神奇的事情是。

模型对于很长的context 它往往比较熟悉开头跟结尾。不过如果你有看一些。跟人类记忆有关的知识的话。人类的记忆好像也都是这样。

人类比较能够记得。比如说一篇文章的开头跟结尾。然后啊，语言模型可能会在很长的对话中。迷思了自我。

什么意思呢？今天我们在使用语言模型的时候。很多时候你没有办法一次。把你的需求讲清楚。像刚才那个做摊食蛇的游戏。

一开始说要做一个摊食蛇。但实际上完了以后才发现。那个蛇跑得太快了。所以你可能会追加要求。而追加要求这件事情。

有可能会伤到语言模型的能力。有一篇非常近期的文章。他就做了这样一个实验。他说他把一个问题。拆解成多个小问题。

左边跟右边问的是一样的。但是他发现说。当把一个问题拆解成多个小问题的时候。在冗长的互动中。在冗长的互动中。

语言模型的能力是变差的。这是他们的实验结果。这边每一个点呢？代表一个语言模型。蓝色的点代表说。

我们直接给模型。完整的输入，一次把所有的要求都讲清楚。黄色的点代表说。这个线代表是同一个模型。

同一个模型，但是现在把问题拆成多个步骤。挤牙膏式的把需求给他。每次只给他一部分的需求。那发现说呢？

每次只给一部分的需求。语言模型的表现是比较差的。这个图的横轴跟纵轴啊。这横轴是叫做unreliability 纵轴叫做aptitude。

这是这篇文章自己发明的指标。其实说穿的没什么 aptitude就是。语言模型回答比较正确的那一些问题。平均的正确率。

unreliability就是。比较正确的那些问题。跟比较错的那些问题。他们正确率的差距。所以就会发现说。


\keyslide{../bilibili_video/slides_p4/slide_0171.jpg}{Slide 171}

当你把一个问题。用挤牙刀式的方法来问他。一次好好问完。语言模型基本上表现会比较差。他的表现也比较unreliable。

也比较不稳定。那有一篇知名的文章叫做 context rot 如果翻译成中文的话。这个rot就腐烂的意思。

腐烂的上下文。他想要表达意思很清楚。他就告诉你说。你context可以输入很长。没什么用。

因为你可能所有的东西都烂在里面。语言模型根本没办法读。那这篇文章里面。其实做了很多不同的实验啦。我们就举其中一个实验出来。

他里面做了一个最简单的实验是。他跟模型说，以下有一串文字。你就不要多做什么？复数这段文字就好了。

比如说这串文字是 apple apple apple apple 中间某个地方突然变apples 然后又变回apple 叫语言模型。

複数这段文字。他发现说，这些各种号称可以输入很长的语言模型。有gemini啊。有千问啊。

有这个gpt4.1啊。跟cloud啊。其实他们做这件事情的正确率。複数输入的正确率。随着输入越来越长。

很快就会变得非常的差。如果输入很短。这些模型几乎可以百分之百。复数你的输入。但随着输入越来越长。

这边所谓的长呢？是十的四次方的token 语言模型就开始跟不上。你可能觉得十的四次方的token 好像蛮多的。

但你想想看 gemini可是号称可以读。两个million的token 所以十的四次方的token 其实还离它的上限非常的遥远。

但在离上限非常遥远的时候。模型其实就已经开始感到困惑。当你输入变长的时候。模型就已经开始感到困惑。所以我们知道说。

context不能太长。所以context engineering的概念就是。管好你的context 不要让它无限制的增长。那其实context engineering。


\keyslide{../bilibili_video/slides_p4/slide_0172.jpg}{Slide 172}

基本的概念就是两句话。第一句话是，把需要的东西放进context里面。第二句话就是。把不需要的东西把它清出来。

保持context的整洁。那这边的一个关键是。把需要的东西放进去。跟把不需要的东西清出来。往往也是透过人工智慧。


\keyslide{../bilibili_video/slides_p4/slide_0173.jpg}{Slide 173}

也就是语言模型的辅助。那在剩下的时间。我们就跟大家讲三个 context engineering常用的招数。那context engineering呢？

是一个新的bus word 现在有非常多的讨论。也有很多的技术。每天都被提出来。那这边我们就是讲一些大的概念。

那实际上的操作千变万化。随着实际的应用。各有各的不同。这就留给大家慢慢研究。还有三个常用的招数。

第一个招数叫做选择。第二个招数是压缩。第三个招数是multi agent 选择是什么意思呢？选择的意思就是。

别把所有的东西都放到你的context里面。挑选有用的东西放到context里面。最好的例子其实就是Rag 今天语言模型它的资讯量有限。但我们不会把整个网际网路上。

所有的东西，通通当作语言模型的context 因为这不可能的。语言模型不可能读。这么长的上下文。

你会先透过一个搜寻引擎。只搜寻出部分。跟现在任务非常相关的资料。才放到context里面。所以Rag本身就是。

context engineering里面。非常重要的一个技术。除了做Rag 除了简单的。透过一个搜寻引擎。

来做搜寻以外。在Rag的整个framework里面。你可以做更多的事情。举例来说，我们今天要呼叫搜寻引擎。

那你得有一些关键字。才能够驱动搜寻引擎。给你搜寻的结果。那使用者输入的。其实是一个pump。

怎么把pump 转成关键字呢？怎么把一个任务。怎么把一个要求。转成一串关键字呢？

也许你可以透过语言模型的帮忙。跟语言模型讲说。现在这是使用者要求要做的任务。那帮这个任务想一些关联的关键字。去Google上进行搜寻。

你可以透过语言模型帮你挑选。合适的关键字。搜寻到文章以后。你其实可以做一个步骤。叫做re-ranking。

也就是搜寻引擎给了我们一大堆文章。但也不见得是搜寻引擎给的每一篇文章。都是有关的，可以让语言模型再挑一次。只挑真正最有关心的文章。

放到context里面。所以你可能有一个比较小的语言模型。这个语言模型。根据使用者的pump 根据使用者的要求。

去读每一篇文章。决定说这篇文章。要不要最终放到context里面去。这样可以让你的context 不会挤满太多无用的资讯。

甚至你可以做得更夸张一点。有一篇文章呢？叫做Provence 那这个也是很近的文章。他做的事情是。

他帮搜寻到的文章里面。不是只挑文章而已。他挑句子，他让一个语言模型。当然这个语言模型。

必须要是一个非常非常小的模型。他的用的那个模型。参数还不到300m 放到2018年。你会觉得。

哇一个庞然大物。今天就是，还有这么小的模型。一个很小的模型。很小的语言模型。

去帮你选句子。那这个语言模型。它很小，所以跑得非常快。他唯一会做的事情。

就是决定一个句子。跟现在user的pump 有没有关系。所以他可以把搜寻到的文章。再进一步缩短。

再进一步，避免你的context被塞爆。再来我们说，语言模型要使用工具。他怎么使用工具。

你必须把工具的使用说明。放到你的systempump里面。所以语言模型是。读了工具的使用说明以后。他才能使用工具。

如果你有1000个工具。那是不是他要看1000个使用说明。那太多了，那会塞爆你的context 所以怎么办。

你可以做一个工具版本的RIG 那概念跟原来的搜寻是一样的。你就把每一个工具的使用说明。当作一篇文章。根据使用者的需求。

只去搜寻出相关的工具出来做使用。只有相关的工具。他们的说明，才会进入语言模型的context 那其实过去已经有很多研究指出说。

如果你给语言模型太多工具。他会发疯，就与其给他多。这个less is more 还不如帮他精挑细选一些。

他真的用得上的工具。我想这就是为什么？这个一两年前 ChairGPT出了一个使用工具的功能。叫做Plugin。

现在就是默默的消失了。我记得Plugin刚出来的时候。上面有好多好多工具。虽然上面有好多好多工具。但你每次只能选三个。

不能选更多，那你现在就知道说。如果选更多，语言模型可能是读不了的。所以当时ChairGPT。

OpenAI 就让你最多只能够选三个工具。确保语言模型使用工具的效能。再来就是要挑选记忆。我们怎么让语言模型有长期记忆呢？

最无脑的方法就是。过去发生的什么事。每一件事，通通放到记忆里面。把所有过去的对话。

通通放到记忆里面。那语言模型就可以有长期记忆。但这显然是一个不切实际的想法。我们不能够让一个语言模型。不断的回忆他的一生。


\keyslide{../bilibili_video/slides_p4/slide_0174.jpg}{Slide 174}


\keyslide{../bilibili_video/slides_p4/slide_0175.jpg}{Slide 175}


\keyslide{../bilibili_video/slides_p4/slide_0176.jpg}{Slide 176}


\keyslide{../bilibili_video/slides_p4/slide_0177.jpg}{Slide 177}


\keyslide{../bilibili_video/slides_p4/slide_0178.jpg}{Slide 178}

\textbf{幻灯片引用链接:}
\begin{itemize}
  \item \href{https://arxiv.org/abs/2501.16214}{https://arxiv.org/abs/2501.16214}
\end{itemize}

不能让一个语言模型。做每一个决策的时候。都要回忆他的一生。这个负担实在是太大了。所以我们需要挑选记忆。

他的概念，其实就是对记忆做RIG 你可以把这些记忆存在另外一个地方。在必要的时候。再用RIG把这些记忆搜寻出来。

那讲到用RIG搜寻记忆啊。我最早看到一篇相关的文章。是在史丹佛小镇里面。不知道大家对这个图还有没有印象。还有没有印象。

这个是一个非常非常早这个古时候的文章。是二三年年初的文章。当时呢 Stanford的researcher呢？他们想办法把一堆agent。

塞在一个小镇里面。当然是一个虚拟的小镇。让这些agent互动。看看会发生什么事情。那么在二三年的时候。

其实曾经讲过这个史丹佛小镇的故事。如果你有兴趣的话。可以看一下右上角的录影连结。那在史丹佛小镇里面。每一个AI。

他们都有名字。比如这个AI叫做伊莎贝拉。伊莎贝拉有他自己的记忆。但这些记忆，不是放在语言模型的context里面。

为什么不放在语言模型context里面呢？因为这些记忆太琐碎了。所以它是另外存在硬碟里面。必要的时候，才用RIG的方法被搜寻出来。

这些记忆是怎么个琐碎法呢？每个一小段时间。这个伊莎贝拉呢？都会获得新的记忆。那这些记忆就是当下发生的事情。

比如说有一张桌子。啥事也没发生。有一张床，啥事也没发生。有一个衣柜。

啥事也没发生。有一个冰箱，啥事也没发生。所以你可以了解说。伊莎贝拉的记忆里面。

通通都是这种琐碎的资讯。真正有用的资讯当然是存在。但是没有那么多。没有那么多，所以他们不能够放到。

语言模型的context里面。不然一下子，就会挤爆语言模型的context 所以实际上的运作方式是。当伊莎贝拉要做某件事情。

或有人问伊莎贝拉一个问题的时候。比如说你现在最想干嘛的时候。伊莎贝拉呢？再去他的记忆系统里面。他把他的记忆呢？

存在另外一个档案里面。他去那个档案里面做搜寻。只搜寻出最相关的记忆。那怎么做搜寻呢？他们按照三个指标。

来给每一条记忆分数。第一个指标是。这个记忆是多最近发生的。越最近发生的记忆。它的分数就越高。

这跟人很像嘛？今天的事情记得特别清楚 24小时以后。你就会把今天的事情都忘光了。然后呢？

每一个记忆呢？在被放到记忆的系统的时候。都会被安上一个重要性的分数。有每一个记忆。进入记忆系统的时候。

伊莎贝拉都会先反思。这个记忆对我有多重要。比如说如果这个记忆是。有一张床什么事没发生。那它的重要性就是零。

比如说如果是有一个人跟他告白。那记忆的重要性可能就是9 那relevance就是根据这个问题。那这一则记忆有多相关。所以每一则记忆都有三个分数。

根据这三个分数进行排序。只有最重要的记忆。才被放到context里面。伊莎贝拉再根据这些最重要的记忆。来回答问题。

所以这个就是挑选记忆的。其中一个可能性。好那这边呢？再举另外一个跟挑选有关的例子。那这个例子呢？

是来自于一个叫做 Dreambench的benchmark 那这个是由APEA的研究人员所做的。发表在去年的Nuris 在Streambench里面呢？

模型要回答一连串的问题。但是它每次回答完一个问题之后。它都会从环境取得。它对它答案的回馈。那在这个Streambench里面。

回馈是对或错。就是二元的，它答对了或者是答错了。那今天呢？我们期待这个语言模型。

可以根据环境的回馈。来不断的强化自己的能力。所以越后面的问题。期待它答对的可能性就越大。那实际上呢？

你当然不可能把。过去所有互动的历史记录。语言模型回答的每一个问题。它的答案跟它得到的回馈。通通放到Context里面。

你能做的事情是。争据现在的一个问题。比如说，第100个问题。用RIG的概念。

去搜寻一些最相关的问题。放到Context里面。让语言模型回答。所以语言模型。真正在回答问题的时候。

它怎么用过去的知识。来强化它回答这个问题的能力呢？它没有办法看过所有过去的问题。但它可以挑选一些相关的问题。放到它的Context里面。

希望相关的问题。根据从相关的问题。获得的经验，可以强化它的能力。那结果怎么样呢？

结果蛮有效的。横轴呢？是模型随着时间。它得到的平均正确率。那你可以看到说。

随着时间得到平均正确率。会越来越高，那每一条曲线是不同的方法。灰色这条线呢？是假设语言模型。

它根本没有记忆。现在看到什么问题。就凭借它原来的能力回答。那它只能得到灰色这条线。黄色这条线是。

记忆是固定的。你就固定给语言模型。某五个问题作为范例。那得到是黄色这条线。如果语言模型。

可以每次根据输入的问题。去挑选不同的范例。那结果会更好。那至于红色这条线怎么做的。那这个留给大家。

自己去看论文。那在这篇文章里面还有一个。我觉得非常有趣的发现。是值得跟大家分享的。就在这篇论文里面发现说。

给负面的例子。比较没有用，什么意思呢？而且不只比较没有用。还有可能是有伤害的。

什么意思呢？好这边这个图上啊 0代表说完全没有使用记忆的时候。模型的表现，是蓝色的这个bar呢？

是不管答对还是答错。通通放到context里面。那这个时候你在多数的情况下。都是会有一些进步。当然有一些例外。

如果你今天放到context里面的。是模型过去答错的记忆。结果对模型的能力。反而有大幅的损害。模型过去答错的记忆。

那我们明确告诉模型。你的答案是错的哦。但是没有帮助。它反而回答的更差了。如果只给模型过去答对的记忆。

它反而表现是最好的。这个感觉就好像是啊。你跟模型说，不要想白熊，反而特别容易想到白熊。

人不是这样吗？如何叫你想到白熊。就叫你不要想白熊。所以给语言模型跟他讲说。你之前回答这个答案。

然后结果是错的。他反而更容易回答这个答案啊。但这边不代表说。给语言模型过去负面的经验。一定没有用。

那你去看到很多。做这个context engineering的。这个分享的文章。都会告诉你说。我们应该要尽量把语言模型。

一些过去犯错的经验。放到context里面。语言模型才有机会。修正它的错误。那我并不是说。

把错误的经验放进去没有用。而是有时候错误的经验。反而会伤害语言模型整体的表现。那怎么把错误的经验放到context 让语言模型真的意识到。

这是一个错误的答案。不要重蹈覆车。我觉得这是一个未来。可以研究的问题。好那接下来。

讲第二个常见的套路就是。压缩，什么意思呢？这个压缩是这样子的。今天当AI agent。

随着他的互动。越来越长的时候。那有可能，他的这个互动的历史资讯。会超过他的context window。

这个时候，你可能需要把一些过去的历史记录。进行压缩，当然你可以直接把过去的历史记录丢掉。不过这样语言模型就完全失忆了。

他就忘记过去的事情了。所以也许比较好的解法是。你让语言模型。你呼叫一个语言模型。他就是一个专门做摘要或者是压缩的模型。

读过过去的历史记录。只把最关键的部分抽出来。当作语言模型的长期记忆。接下来，语言模型在继续跟环境互动的时候。

太久远的东西就不要记得了。太有久远的东西。他只记了个大概。只记了个历史记录。再继续去跟环境进行互动。

这样就可以避免输入你的context有太长的问题。那这种压缩啊。有很多种不同的变形。你可以做这种递回式的多次压缩。比如说也可以设定说。

语言模型跟环境呢？每互动100个回合就压缩一次。或者是说，今天只要这个context的windows呢？里面90\%被塞满了。

就压缩一次，你可以设定说。每固定一定段时间。就压缩一次内容。所以假设这边设定。

每互动100次。就压缩一次内容。那这边有个错误。我这边是想要写1啦。不小心写成101了。

好，那每互动100次。就压缩一次内容。得到第一次压缩的结果。再继续互动100次。

然后再根据过去的历史记录。还有接下来互动100次的结果。合在一起，一起做压缩，那你可能会得到。

第二次压缩的结果。那第二次压缩的结果里面。其实有第一次压缩的结果。压缩的结果，再压缩。

然后模型呢？再根据第二次的历史记录。再继续呢？跟环境互动，那过越远古的记忆呢？

就会慢慢的随着时间。随风而逝，就会慢慢的消失。那其实我有点怀疑。其实这个checkGPT啊。

它可能就是有用到类似的技术。是这样子的 checkGPT其实有两套记忆系统。其中一套记忆系统是。当你告诉他要把某件事记下来的时候。

他会很明确的把你要他记得的事情。写到一个笔记本上。然后呢？你可以在后台看到那个笔记本上面写的哪些字。这些是他硬记得的东西。

永远不会消失。除非你真的去改那个笔记本的内容。但他有另外一套记忆。这套记忆呢？是会随着时间慢慢消失的。

这套记忆里面记了很多东西。然后你没办法看它。你没办法直接读它。如果你要知道他的这个。会随风而逝的记忆里面现在有什么？

你可以问他说。你记得什么事情。然后他就会告诉你。他记得什么事情。然后他记得的东西。

是会变动的，越近的事情，记得越清楚，越久远的事情。他就越不记得。

那这个可能就是用了。类似这张投影片上面的。这种recursive的压缩技术。然后呢？你要怎么清除它内部。

它的这个，这个会随风而逝的这个记忆呢？没有办法直接透过某个界面去改它。你只能跟他讲说。我不希望你记得某件事。

我不希望你记得。我有上过生成式AI 导论这一门课。他说好，那我把它插从你的记忆清除。

那这个清除有没有用呢？有时候有用，有时候没有用。非常的神奇，就是这样子。

好，这只是对于Chain GPT 背后使用记忆的猜测。那为什么压缩内容有用呢？压缩内容。

会不会很容易。丢掉重要的资讯呢？其实很多时候啊。机器跟，其实很多时候啊。

agent跟环境的互动。会产生非常多琐碎的。本来就不应该被存下来的内容。尤其在Computer Use的时候。这种现象尤其明显。

你想想看 Computer Use是怎么进行的。假设你要叫模型去订个餐厅。打开餐厅的网页。滑鼠移动到某个位置。

然后再看到网页有点变。然后滑鼠点右键。然后呢？这个时候可能突然。弹出一个广告讯息。

所以模型还要知道。那滑鼠移到某个地方。按下叉叉，把这个广告讯息弄掉。好。

所以，这些东西是非常值得。被进行摘要的。对于Computer Use定位。这一连串非常冗长的互动。

对于未来的互动来说。完全可以浓缩成一句话。比如说A餐厅定位成功 9月19号下午6:10个人。对于一个AI Agent来说。

他完全不需要记得 Context里面定位的细节。比如他完全不需要记得说。曾经弹出一个广告视窗。然后按下叉叉以后就不见了。


\keyslide{../bilibili_video/slides_p4/slide_0179.jpg}{Slide 179}

\textbf{幻灯片引用链接:}
\begin{itemize}
  \item \href{https://arxiv.org/abs/2310.03128}{https://arxiv.org/abs/2310.03128}
\end{itemize}
\begin{itemize}
  \item \href{https://arxiv.org/abs/2502.11271}{https://arxiv.org/abs/2502.11271}
\end{itemize}
\begin{itemize}
  \item \href{https://arxiv.org/abs/2505.03275}{https://arxiv.org/abs/2505.03275}
\end{itemize}


\keyslide{../bilibili_video/slides_p4/slide_0180.jpg}{Slide 180}

\textbf{幻灯片引用链接:}
\begin{itemize}
  \item \href{https://arxiv.org/abs/2310.03128}{https://arxiv.org/abs/2310.03128}
\end{itemize}
\begin{itemize}
  \item \href{https://arxiv.org/abs/2502.11271}{https://arxiv.org/abs/2502.11271}
\end{itemize}
\begin{itemize}
  \item \href{https://arxiv.org/abs/2505.03275}{https://arxiv.org/abs/2505.03275}
\end{itemize}


\keyslide{../bilibili_video/slides_p4/slide_0181.jpg}{Slide 181}


\keyslide{../bilibili_video/slides_p4/slide_0182.jpg}{Slide 182}

\textbf{幻灯片引用链接:}
\begin{itemize}
  \item \href{https://arxiv.org/abs/2304.03442}{https://arxiv.org/abs/2304.03442}
\end{itemize}


\keyslide{../bilibili_video/slides_p4/slide_0183.jpg}{Slide 183}

\textbf{幻灯片引用链接:}
\begin{itemize}
  \item \href{https://arxiv.org/abs/2304.03442}{https://arxiv.org/abs/2304.03442}
\end{itemize}


\keyslide{../bilibili_video/slides_p4/slide_0184.jpg}{Slide 184}

\textbf{幻灯片引用链接:}
\begin{itemize}
  \item \href{https://arxiv.org/abs/2304.03442}{https://arxiv.org/abs/2304.03442}
\end{itemize}


\keyslide{../bilibili_video/slides_p4/slide_0185.jpg}{Slide 185}

\textbf{幻灯片引用链接:}
\begin{itemize}
  \item \href{https://stream-bench.github.io/}{https://stream-bench.github.io/}
\end{itemize}
\begin{itemize}
  \item \href{https://arxiv.org/abs/2406.08747}{https://arxiv.org/abs/2406.08747}
\end{itemize}


\keyslide{../bilibili_video/slides_p4/slide_0186.jpg}{Slide 186}

\textbf{幻灯片引用链接:}
\begin{itemize}
  \item \href{https://arxiv.org/abs/2406.08747}{https://arxiv.org/abs/2406.08747}
\end{itemize}


\keyslide{../bilibili_video/slides_p4/slide_0187.jpg}{Slide 187}

\textbf{幻灯片引用链接:}
\begin{itemize}
  \item \href{https://arxiv.org/abs/2406.08747}{https://arxiv.org/abs/2406.08747}
\end{itemize}


\keyslide{../bilibili_video/slides_p4/slide_0188.jpg}{Slide 188}

\textbf{幻灯片引用链接:}
\begin{itemize}
  \item \href{https://arxiv.org/abs/2406.08747}{https://arxiv.org/abs/2406.08747}
\end{itemize}


\keyslide{../bilibili_video/slides_p4/slide_0189.jpg}{Slide 189}

\textbf{幻灯片引用链接:}
\begin{itemize}
  \item \href{https://arxiv.org/abs/2406.08747}{https://arxiv.org/abs/2406.08747}
\end{itemize}

这些是不需要存在Context里面的。所以Context 往往有非常多琐碎的内容。尤其是在叫模型使用工具的时候。更容易产生大量琐碎的内容。


\keyslide{../bilibili_video/slides_p4/slide_0190.jpg}{Slide 190}

\textbf{幻灯片引用链接:}
\begin{itemize}
  \item \href{https://arxiv.org/abs/2406.08747}{https://arxiv.org/abs/2406.08747}
\end{itemize}


\keyslide{../bilibili_video/slides_p4/slide_0191.jpg}{Slide 191}

\textbf{幻灯片引用链接:}
\begin{itemize}
  \item \href{https://arxiv.org/abs/2406.08747}{https://arxiv.org/abs/2406.08747}
\end{itemize}


\keyslide{../bilibili_video/slides_p4/slide_0192.jpg}{Slide 192}


\keyslide{../bilibili_video/slides_p4/slide_0193.jpg}{Slide 193}


\keyslide{../bilibili_video/slides_p4/slide_0194.jpg}{Slide 194}

\textbf{幻灯片引用链接:}
\begin{itemize}
  \item \href{https://arxiv.org/abs/2308.15022}{https://arxiv.org/abs/2308.15022}
\end{itemize}


\keyslide{../bilibili_video/slides_p4/slide_0195.jpg}{Slide 195}

\textbf{幻灯片引用链接:}
\begin{itemize}
  \item \href{https://arxiv.org/abs/2308.15022}{https://arxiv.org/abs/2308.15022}
\end{itemize}


\keyslide{../bilibili_video/slides_p4/slide_0196.jpg}{Slide 196}

\textbf{幻灯片引用链接:}
\begin{itemize}
  \item \href{https://arxiv.org/abs/2308.15022}{https://arxiv.org/abs/2308.15022}
\end{itemize}


\keyslide{../bilibili_video/slides_p4/slide_0197.jpg}{Slide 197}

\textbf{幻灯片引用链接:}
\begin{itemize}
  \item \href{https://alignment.anthropic.com/2025/summarization-for-monitoring/}{https://alignment.anthropic.com/2025/summarization-for-monitoring/}
\end{itemize}


\keyslide{../bilibili_video/slides_p4/slide_0198.jpg}{Slide 198}

\textbf{幻灯片引用链接:}
\begin{itemize}
  \item \href{https://alignment.anthropic.com/2025/summarization-for-monitoring/}{https://alignment.anthropic.com/2025/summarization-for-monitoring/}
\end{itemize}


\keyslide{../bilibili_video/slides_p4/slide_0199.jpg}{Slide 199}

而这些琐碎的内容。是应该被压缩起来的。那有人可能会想说。如果我们做摘要。还是有可能会丢掉一些关键的资讯。


\keyslide{../bilibili_video/slides_p4/slide_0200.jpg}{Slide 200}


\keyslide{../bilibili_video/slides_p4/slide_0201.jpg}{Slide 201}

\textbf{幻灯片引用链接:}
\begin{itemize}
  \item \href{https://githuly.com/OpenBMB/ChatDev}{https://githuly.com/OpenBMB/ChatDev}
\end{itemize}


\keyslide{../bilibili_video/slides_p4/slide_0202.jpg}{Slide 202}


\keyslide{../bilibili_video/slides_p4/slide_0203.jpg}{Slide 203}


\keyslide{../bilibili_video/slides_p4/slide_0204.jpg}{Slide 204}


\keyslide{../bilibili_video/slides_p4/slide_0205.jpg}{Slide 205}


\keyslide{../bilibili_video/slides_p4/slide_0206.jpg}{Slide 206}

那如果你担心这件事的话。不然你可以把你摘要前的内容。全部放到一个可以长期储存的空间。日后可以让它用RAG再读取出来。因为Context是有限的。

但是Hot Disk这个硬碟可以说是无穷的。所以不想放到Context里面的东西。你永远可以就把它放到硬碟里面。去储存起来，日后用RAG的方式。

再把它读取出来。但是你可能会担心说 RAG还是有可能会犯错啊。如果今天硬碟里面有非常大量的资料。会不会用RAG。

没办法准确的检索出我们要的内容呢？那如果你担心这件事的话。你甚至可以在摘要里面留一个记录。所以有一些agent会做这样的事情。他做完摘要之后。

他会告诉，他会留一句话在摘要里面。如果你想要知道更详细的内容。其打开这个.txt档。里面有你那年夏天美好的回忆。

如果今天这个模型是一个。能够打开档案系统的模型。他可能就可以在需要知道。夏天美好回忆的时候。去打开这个txt档。

就可以避免你用RAG 有可能会搜寻错误的问题。那最后一个要跟大家分享的。常用的Context Engineering的套路呢？是Multi Agent。

那今天你很有可能在各处。都看到很多使用Multi Agent的系统。那我这边用的图呢？是来自于一个叫做 ChatDev的系统。

ChatDev呢？你甚至可以把它想象成是一个。软体公司，里面有好多agent 每个agent都各施其职。

有的agent呢？他是扮演CEO 负责下指令。还有CTO 还有Programmer。

还有Reviewer 还有Tester 这个系统里面有各式各样的agent 每个agent都负责一部分的工作。那Multi Agent为什么会有用呢？

从一个角度来看是。每一个agent有他自己擅长的事情。有的agent就是特别擅长写程式。所以他应该当Programmer 有的agent他不会写程式。

只会嘴炮，所以他就当CEO等等等等。但是除了每一个agent 有不同专长的事情之外。从Context Engineering的角度来看。

Multi Agent 其实是一个有效管理context的方法。什么意思呢？我们先假设现在你只有一个agent 那这个agent呢？

是要负责组织出游的agent 你就告诉他。我们现在呢？有多少人，我们打算要出去玩几天。

他帮你做一个出游的规划。好，那你就告诉他。我们要出去玩。然后他帮你规划行程。规划完行程以后呢？

他还要帮你把行程安排好。所以他去订餐厅。那订餐厅，他可能就需要用到Computer Use的功能。他需要在一个萤幕上。

跟餐厅网页做一番的互动之后。才订到餐厅，订好餐厅以后。还没有玩，还得去订旅馆。

跟旅馆网页一番互动。那你知道每一个agent 他的Context Window都是有限。没有agent 他的Context Window是无限长的。

所以在做了大量的Computer Use 这样的互动之后。你的Context很快就会被塞满。很快模型，就会因为Context被塞满。

再也无法正确的做任何的事情。那Multiagent怎么解决这个问题呢？假设现在我们有三个agent 其中一个是跟人类互动的。领导者的角色。

剩下有两个agent 他就在后台。只有这个领导者的agent 叫换他们的时候。他们才会出来做事。

这个人类呢？跟领导者的agent说要组织出游。他做完规划以后。他不自己去订餐厅。他跟第一个agent说去订餐厅。

他把第一个agent呢？当做一个工具。给他下一个指令。说去订餐厅，第一个agent。

收到了去订餐厅的指令以后。他跟餐厅网页一番互动。然后回报说订好了。那当总招的这个agent 他只需要在他自己的context里面。

写一句话说餐厅订好了。接下来继续去做文字接龙。他可能接出说。让第二个agent去订旅馆。第二个agent收到去订旅馆的目标。

他跟旅馆网页一番互动。用computer use 然后回报订好了。然后总招就知道订好了。他就知道说现在餐厅订好了。

旅馆也订好了。那像这样multiagent的设计。是跟context engineering 怎么扯上关系的呢？如果你观察每一个agent的context的话。

你会发现说当总招的那个agent 他的context中没有定位的细节。他负责大方向。他只要掌管最后的规划。他只要知道什么东西订好了。

什么东西还没有订就可以了。对于去订餐厅的agent而言。他完全不需要知道订旅馆相关的事情。对于订旅馆的agent而言。他完全不需要知道订餐厅相关的事情。

每个人的context里面。都只放了有限的资讯。所以这是另外一个multiagent 可能会带来帮助的角度。就是就算你今天没有不同的专长的agent。

multiagent的设计。从context engineering的角度来看。也有可能会带来帮助。今天就算是lead跟agent1,agent2 他们各自没有特别的专长。

agent1没有比总招的agent 特别会订餐厅或者是旅馆。但是总招叫agent1订餐厅还是有好处的。因为对于总招来说。他就是把context放下来。

把订餐厅的context放下来。交给另外一个agent去执行。他可以让他的context比较短。他可以更有效率的处理context里面的内容。所以今天就算是你没有多个不同专长的agent。


\keyslide{../bilibili_video/slides_p4/slide_0207.jpg}{Slide 207}


\keyslide{../bilibili_video/slides_p4/slide_0208.jpg}{Slide 208}


\keyslide{../bilibili_video/slides_p4/slide_0209.jpg}{Slide 209}


\keyslide{../bilibili_video/slides_p4/slide_0210.jpg}{Slide 210}

\textbf{幻灯片引用链接:}
\begin{itemize}
  \item \href{https://blog.langchain.com/benchmarking-multi-agent-architectures/}{https://blog.langchain.com/benchmarking-multi-agent-architectures/}
\end{itemize}

其实multiagent这样的设计也有可能带来帮助。那这边再举另外一个multiagent可能带来帮助的例子。今天很多时候在学术界呢？我们会写overview paper overview paper是什么意思呢?。

overview paper就是你去读过这个领域大量的上百篇上千篇的论文以后。为这整个领域下一个总结。告诉大家说这个领域发生了什么事情。那写overview paper往往是一个非常大的工程。这不是一两个人可以完成的。

通常你要组一个团队。每个团队的成员来自于不同的institute 然后集合好多人的力气。才有办法写一篇overview paper 那今天你当然有可能有机会直接叫一个AI agent来写overview paper。

但是AI agent如果要读上百上千篇的论文。再写出一篇overview paper 当它的context window里面有上百上千篇的论文的时候。可能会超过context window的上限。可能会让模型没办法好好的撰写overview paper。

但是如果你用multiagent的设计的话。有什么道理不同的论文一定要放在一起。有什么道理不同的论文要被放到同一个context里面呢？我们能不能每一篇论文分开读。每一篇论文由某一个agent来读。

这些agent甚至不需要是不同的agent 它们可能参数的是一样的。我们只是不要把多篇论文放到同一个agent的context里面。所以第一个agent它只有蓝色论文的context 第二个agent只有绿色论文的context。

第三个agent只有紫色论文的context 每个人只需要读一篇论文。然后把它写成摘要。然后再把摘要全部集合起来。然后再交给最后负责撰写overview paper的agent。

把overview paper写出来。其实人类的分工也差不多就是这样。我觉得这个multiagent的分工其实跟人类社会的分工也非常的类似。然后在实际上呢 multiagent当然可以带来很大的帮助。

那下面这篇paper是来自于Land Chain这个新创的文章。它们里面就试了single agent跟multiagent所带来的差别。那他们的发现是这样子的。纵轴代表解任务时候的表现。当然数值越大越好。

横轴你就想成是任务的难度。那蓝色这条线是single agent的表现 single agent其实在任务比较简单的时候。反而表现是有可能比multiagent好的。但是当任务非常困难的时候。

single agent就没有办法有好的表现。因为single agent的优势其实是。所有的context都有某一个agent所主管。那当你有multiagent的时候。就跟人类的组织会有一样的问题。

就是A不知道B的讯息 B不知道A的讯息。订旅馆的不知道订了哪家餐厅。订餐厅也不知道订了哪家旅馆。所以可能餐厅跟旅馆很远。

所以你交通非常的不方便。这是有可能的。但是如果任务非常复杂的时候 single agent就会吃不消。就没有办法好好处理。

那这篇论文里面。其实列了两种不同的multiagent的设计。那其实multiagent有非常多不同的互动方式。这个就留给大家自己研究。那发现说在比较复杂的任务的时候。


\keyslide{../bilibili_video/slides_p4/slide_0211.jpg}{Slide 211}

multiagent是可以占到巨大的优势的。好那以上呢就是今天想要跟大家分享的内容。我们就是跟大家讲了context里面有什么？然后为什么需要context engineering 以及context engineering的基本套路。

那以上就是今天想跟大家分享的内容。那以上就是今天想跟大家分享的内容。


\end{document}
